\documentclass[../../main.tex]{subfiles}
\begin{document}

\chapter{The Transformer architecture and attention mechanisms}
\label{ch:transformer_architecture}

Chapter~\ref{ch:word_embeddings} described how tokenization and embedding layers map discrete input sequences to vectors with positional information. Chapter~\ref{ch:rnn_sequence_modeling} showed how recurrent networks process these vectors by updating a hidden state over time.

Recurrence imposes a sequential dependency: computing $\vh_t = f(\vh_{t-1}, \vx_t)$ requires the state from the preceding position, which limits parallel computation across a sequence. In encoder-decoder networks, compressing the input into a fixed-size state can also restrict the information available to the decoder. Attention addresses the latter limitation by allowing representations at one position to use information from other positions; the Transformer applies attention without a recurrent state update.

Vaswani et al. introduced the Transformer in 2017, using attention in place of recurrence \cite{vaswani2017}. This chapter derives scaled dot-product attention and examines bidirectional encoder self-attention, causal decoder self-attention, and cross-attention. It then develops multi-head projections, residual connections, Pre-LN and Post-LN normalization, and relative positional representations. The discussion also compares encoder-decoder models such as T5, encoder-only models such as BERT, and decoder-only autoregressive models, before treating sampling methods for sequence generation.

\clearpage
\section{The Self-Attention Primitive}
\label{sec:recurrent_impasse}

For nearly a decade, sequence modeling was dominated by Recurrent Neural Networks (RNNs) and their gated variants, Long Short-Term Memory (LSTM) networks and Gated Recurrent Units (GRUs). As established in Chapter~\ref{ch:rnn_sequence_modeling}, recurrent models represent a sequence of tokens by recursively updating an internal hidden state vector:
\begin{equation}
\label{eq:recurrent_transition_step}
\vh_t = f(\vh_{t-1}, \vx_t), \quad t \in \{1, 2, \dots, L\},
\end{equation}
where $\vh_t \in \R^{d_h}$ encapsulates the historical context observed up to time step $t$, and $\vx_t \in \R^d$ is the input vector representation at time $t$.

While gating mechanisms and Constant Error Carousels mitigated the vanishing gradient problem over moderate temporal spans, recurrence inherently imposes two severe, prohibitive bottlenecks that halted further scaling in language modeling:

\subsubsection{The sequential execution wall}
The transition equation $\vh_t = f(\vh_{t-1}, \vx_t)$ enforces an unbroken, strictly sequential causal dependency chain: computation of $\vh_t$ cannot commence until the physical memory write of $\vh_{t-1}$ is completed. Consequently, an input sequence of $L$ tokens mandates $\mathcal{O}(L)$ sequential computational operations.
Modern deep learning accelerators (GPUs and TPUs) derive their performance from massive data parallelism, executing thousands of arithmetic matrix operations concurrently across high-bandwidth tensor cores. Under recurrent processing, hardware parallelism across the temporal dimension is impossible: the accelerator is forced to idle between consecutive time steps, creating a computational wall that restricts training across massive datasets.

\subsubsection{The information compression bottleneck of Seq2Seq encoders}
In classical sequence-to-sequence architectures, an encoder compresses an input sequence of arbitrary length $L$ into a single, fixed-size vector $\vh_L \in \R^{d_h}$ (or a tuple $(\vh_L, \vc_L)$ in LSTMs). This single vector is subsequently handed over to the decoder to initiate target sequence generation.
As the input length $L$ expands---such as in multi-turn dialogues, technical documentation, or software engineering source code files spanning thousands of tokens---compressing all semantic, syntactic, and structural relationships into a finite-dimensional vector becomes physically impossible. Information from early time steps undergoes continuous dissipation, inducing severe contextual amnesia.


These limitations raised a fundamental question: \emph{Can we design a neural architecture that completely eliminates recurrence, operates across all tokens in parallel with $\mathcal{O}(1)$ sequential operations, and allows every token to directly communicate with every other token in the sequence without a fixed-size information bottleneck?}
In 2017, this question was answered by the introduction of the \textbf{Transformer} \cite{vaswani2017}, an architecture built entirely upon multi-head attention mechanisms.

\clearpage
\section{Transformer Architecture Overview}
\label{sec:high_level_encoder_decoder}

The foundational Transformer proposed in \cite{vaswani2017} adopts an \textbf{Encoder-Decoder} paradigm, structured to eliminate the recurrent compression bottleneck while enabling end-to-end parallel training. Figure~\ref{fig:vaswani_original_architecture} reproduces the canonical architectural diagram from the landmark paper \cite{vaswani2017}.

\begin{figure}[htbp]
\centering
\includegraphics[width=0.52\textwidth]{figures/part1/transformer_architecture_vaswani2017.png}
\caption{The original Transformer model architecture as published in \cite{vaswani2017}. The left stack corresponds to the Encoder ($N=6$ identical layers), and the right stack corresponds to the Decoder ($N=6$ identical layers), equipped with causal masked self-attention and cross-attention over the encoder representations.}
\label{fig:vaswani_original_architecture}
\end{figure}

\subsubsection{The Encoder stack: uncompressed contextual representations}
The Encoder consists of a stack of $N$ identical layers (typically $N=6$ in the base model). Crucially, the Encoder does \emph{not} compress the input sequence into a single summary vector. Instead, it takes the input sequence of position-augmented embeddings $\mX \in \R^{L_{\text{enc}} \times \dmodel}$ and transforms it into an equally sized sequence of contextual representations:
\begin{equation}
\label{eq:encoder_uncompressed_output}
\mH_{\text{enc}} = [\vh_1, \vh_2, \dots, \vh_{L_{\text{enc}}}] \in \R^{L_{\text{enc}} \times \dmodel}.
\end{equation}
Every output vector $\vh_i$ possesses a global receptive field that incorporates information from every other token in the sequence through bidirectional self-attention, preserving fine-grained token-level semantics without information loss.

\subsubsection{The Decoder stack: autoregressive generation and cross-attention}
The Decoder also consists of a stack of $N$ identical layers, operating autoregressively at inference time:
\begin{enumerate}
    \item \textbf{Autoregressive sequence initialization}: Generation is primed by feeding the special beginning-of-sequence delimiter $\langle\text{BOS}\rangle$ as the initial token.
    \item \textbf{Causal step-by-step prediction}: The Decoder processes the tokens generated so far, attends to the encoder representations $\mH_{\text{enc}}$, and emits logits over vocabulary $\vocab$. The highest-scoring token (or sampled token) $y_1$ is generated.
    \item \textbf{Prefix feedback loop}: The newly emitted token is appended to the decoder input, forming the context $[\langle\text{BOS}\rangle, y_1]$ to predict $y_2$. This loop repeats until the model outputs the end-of-sequence delimiter $\langle\text{EOS}\rangle$, terminating generation.
\end{enumerate}

\subsubsection{The training phase: Teacher Forcing and sequence parallelization}
At inference, the decoder must generate tokens sequentially. However, during training, the entire target sequence $[y_0, y_1, \dots, y_{L_{\text{dec}}-1}]$ is already known.
By employing \textbf{Teacher Forcing} combined with causal attention masking (which prevents tokens from attending to subsequent target tokens), all $L_{\text{dec}}$ generation steps are computed simultaneously in a single forward pass. The decoder does not unroll an iterative loop during training: the loss across all output positions is evaluated concurrently on the GPU.

\clearpage
\section{Attention Principles}
\label{sec:attention_first_principles}

\subsection{Why we need the attention mechanism: from RNN bottlenecks to dynamic interaction}
\label{subsec:why_we_need_attention}

To understand the conceptual necessity of the attention mechanism, we must examine the architectural bottleneck that afflicted sequential models prior to the Transformer, as formalized by Sebastian Raschka \cite{raschka2024llm}.

When translating natural language text (e.g., from German or Italian to English) or performing code generation from formal specifications, literal word-by-word substitution is fundamentally impossible. Grammatical structures, morphological agreements, and clause ordering diverge substantially between linguistic systems. Consequently, sequence transduction requires a deep model consisting of two interacting submodules: an \textbf{Encoder} that ingests and understands the source sequence, and a \textbf{Decoder} that synthesizes the target sequence.

Prior to Transformers, the dominant paradigm for sequence-to-sequence modeling was the recurrent neural network (RNN, Chapter~\ref{ch:rnn_sequence_modeling}). In an encoder-decoder RNN:
\begin{enumerate}
    \item The encoder ingests source tokens $\vx_1, \vx_2, \dots, \vx_L$ sequentially, updating its recurrent hidden state at each time step via $\vh_t = f(\vh_{t-1}, \vx_t)$.
    \item At the conclusion of the source sequence, the encoder attempts to compress the entire semantic and syntactic content of the input into its final hidden state vector $\vh_L \in \R^{d_h}$.
    \item This single vector $\vh_L$ is handed over to the decoder as its initial context vector $\vc = \vh_L$, from which the decoder must generate the translated output token by token.
\end{enumerate}

This architecture suffers from a substantial mathematical limitation: the \textbf{information compression bottleneck} \cite{raschka2024llm}. Forcing an arbitrary-length sentence—spanning tens, hundreds, or thousands of tokens—into a single fixed-dimensional embedding vector creates an information-theoretic bottleneck. As sequence length expands:
\begin{itemize}
    \item The decoder possesses \emph{zero direct access} to intermediate representations $\vh_1, \vh_2, \dots, \vh_{L-1}$ from earlier encoder steps.
    \item Gradients backpropagating through time undergo exponential decay or explosion across recurrent chain products $\prod_{k=t+1}^L \frac{\partial \vh_k}{\partial \vh_{k-1}}$.
    \item Information from early tokens undergoes continuous dissipative dilution, causing severe context amnesia on long sequences.
\end{itemize}

To solve this dilemma, Bahdanau et al.~\cite{bahdanau2014} introduced the attention mechanism for RNNs, providing the decoder with dynamic, direct connections to \emph{all} intermediate encoder states $\vh_1, \dots, \vh_L$ at every generation step.
Three years later, Vaswani et al.~\cite{vaswani2017} realized that recurrence itself was entirely superfluous. By introducing \textbf{Self-Attention}, the model allows every position within a sequence to directly attend to and dynamically interact with every other position in the sequence across an $\mathcal{O}(1)$ sequential path length, computing rich contextual embeddings where each word's representation is dynamically updated based on the relevant context of the entire sentence \cite{raschka2024llm}.

\subsection{The associative memory analogy: queries, keys, and values}
\label{subsec:associative_memory_analogy}

To understand how attention operates mathematically from first principles, it is instructive to draw an analogy to associative retrieval structures in computer science \cite{jurafsky2024, sanderson2024attention}.

\subsubsection{Classical dictionary lookup (discrete hash table)}
In computer systems, a dictionary or associative array stores information as a collection of $m$ discrete key-value pairs:
\begin{equation}
\mathcal{D} = \{ (\kappa_1, \nu_1), (\kappa_2, \nu_2), \dots, (\kappa_m, \nu_m) \}.
\end{equation}
When a discrete search query $q$ is issued, the dictionary performs an exact binary match against the keys:
\begin{equation}
\label{eq:discrete_lookup}
\text{\texttt{Lookup}}(q, \mathcal{D}) = \begin{cases}
\nu_k & \text{if } \exists\, k \text{ such that } q = \kappa_k, \\
\text{\texttt{None}} & \text{otherwise}.
\end{cases}
\end{equation}
In this classical paradigm, retrieval is binary ($0$ or $1$) and non-differentiable: if query $q$ is slightly perturbed, the lookup either matches completely or fails substantially.

\subsubsection{Soft, differentiable key-value retrieval}
Attention generalizes dictionary lookup to continuous metric vector spaces $\R^d$:
\begin{itemize}
    \item A \textbf{Query} vector $\vq \in \R^{d_k}$ represents the information need: \emph{``What linguistic, syntactic, or semantic attributes am I searching for?''}
    \item A set of $m$ \textbf{Key} vectors $\vk_j \in \R^{d_k}$ represent content identifiers or addresses: \emph{``What attributes do I offer?''}
    \item A set of $m$ \textbf{Value} vectors $\vv_j \in \R^{d_v}$ represent the informative content or payload: \emph{``What information do I communicate if selected?''}
\end{itemize}

Instead of checking exact equality $q = \kappa_k$, we evaluate a continuous compatibility score $s(\vq, \vk_j) \in \R$ using vector inner products:
\begin{equation}
\label{eq:raw_compatibility_score}
s_j = \vq^\top \vk_j.
\end{equation}
The raw alignment scores $[s_1, s_2, \dots, s_m]$ are mapped onto the probability simplex via the $\softmax$ operator, yielding non-negative attention weights that sum to unity:
\begin{equation}
\label{eq:softmax_attention_weight}
\alpha_j = \frac{\exp(s_j)}{\sum_{l=1}^m \exp(s_l)}, \quad \sum_{j=1}^m \alpha_j = 1, \quad \alpha_j \ge 0.
\end{equation}
Finally, the retrieved representation is computed as the \textbf{convex combination} (weighted sum) of all value vectors:
\begin{equation}
\label{eq:attention_convex_combination}
\vo = \sum_{j=1}^m \alpha_j \vv_j \in \R^{d_v}.
\end{equation}
Because inner products, exponentials, and summations are smooth operations, the entire lookup pipeline is fully differentiable with well-defined gradients $\frac{\partial \vo}{\partial \vq}$, $\frac{\partial \vo}{\partial \vk_j}$, and $\frac{\partial \vo}{\partial \vv_j}$.

\subsubsection{Illustrative example: contextual displacement of a fluffy blue creature}
To understand how soft key-value retrieval transforms static representations into rich contextual embeddings in high-dimensional vector spaces, consider the visual intuition formalized by Grant Sanderson \cite{sanderson2024attention} and illustrated in Figure~\ref{fig:3b1b_query_key_creature}:
\begin{enumerate}
    \item \textbf{The static embedding}: Initially, each word in the phrase \emph{``A fluffy blue creature roamed the verdant forest''} begins as an isolated unigram embedding vector $\ve \in \R^{\dmodel}$. The static vector for \emph{``creature''} merely occupies a generic semantic coordinate indicating an animal or organism, completely devoid of context regarding its visual appearance.
    \item \textbf{Queries as informational requests}: When passed into an attention block, linear projection $\mW_Q$ maps the vector for \emph{``creature''} into a Query vector $\vq_{\text{creature}}$. Geometrically, this query acts as a directional probe asking: \emph{``What descriptive adjectives precede me?''}
    \item \textbf{Keys as feature advertisements}: Concurrently, projection matrix $\mW_K$ maps all tokens into Key vectors. The preceding adjectives \emph{``fluffy''} and \emph{``blue''} produce keys $\vk_{\text{fluffy}}$ and $\vk_{\text{blue}}$ advertising: \emph{``I am a tactile descriptor''} and \emph{``I am a visual color descriptor''}.
    \item \textbf{Inner-product compatibility}: Computing inner products $\vq_{\text{creature}}^\top \vk_j$ evaluates geometric alignment (Figure~\ref{fig:3b1b_attn_scores}). Because the query explicitly seeks adjective descriptors, it achieves high positive inner products with $\vk_{\text{fluffy}}$ ($+93.0$) and $\vk_{\text{blue}}$ ($+93.4$), but low or negative alignment with irrelevant tokens like \emph{``roamed''} or \emph{``forest''}.
    \item \textbf{Value routing and displacement}: Linear projection $\mW_V$ extracts the informative payload: Value vectors $\vv_{\text{fluffy}}$ and $\vv_{\text{blue}}$ encode specific sensory features. Taking their weighted sum produces a contextual displacement vector (Figure~\ref{fig:3b1b_value_table}):
    \begin{equation}
    \label{eq:fluffy_blue_displacement}
    \Delta \ve_{\text{creature}} = \alpha_{\text{fluffy}} \vv_{\text{fluffy}} + \alpha_{\text{blue}} \vv_{\text{blue}}.
    \end{equation}
    When added to the initial representation via residual connection, $\ve_{\text{creature}} + \Delta \ve_{\text{creature}}$ physically moves the point in embedding space into the specialized manifold region corresponding to a \emph{fluffy, blue creature}.
\end{enumerate}

\begin{figure}[htbp]
\centering
\begin{subfigure}[b]{0.48\textwidth}
    \centering
    \resizebox{\textwidth}{!}{%
    \begin{tikzpicture}[
        cell/.style={rectangle, draw=lightGrayLine, minimum size=7.2mm, thick, font=\tiny\sffamily, align=center},
        hcell/.style={rectangle, draw=examAmberLine, fill=examAmberBg, minimum size=7.2mm, thick, font=\tiny\bfseries\sffamily, align=center},
        qlabel/.style={font=\tiny\bfseries\sffamily, align=right, anchor=east},
        klabel/.style={font=\tiny\bfseries\sffamily, align=center, rotate=45, anchor=south west}
    ]
        % Column headers (Keys)
        \foreach \t [count=\j from 0] in {A, fluffy, blue, creature, roamed, forest} {
            \node[klabel] at (\j*0.8 + 0.4, 0.2) {\t};
        }
        % Row 0: A
        \node[qlabel] at (-0.1, -0.4) {A};
        \node[cell, fill=white] at (0.4, -0.4) {+1.2};
        \node[cell, fill=white] at (1.2, -0.4) {-0.8};
        \node[cell, fill=white] at (2.0, -0.4) {-1.1};
        \node[cell, fill=white] at (2.8, -0.4) {+4.3};
        \node[cell, fill=white] at (3.6, -0.4) {-0.2};
        \node[cell, fill=white] at (4.4, -0.4) {+0.5};
        
        % Row 1: fluffy
        \node[qlabel] at (-0.1, -1.2) {fluffy};
        \node[cell, fill=white] at (0.4, -1.2) {-0.5};
        \node[cell, fill=white] at (1.2, -1.2) {+3.1};
        \node[cell, fill=white] at (2.0, -1.2) {+2.8};
        \node[cell, fill=white] at (2.8, -1.2) {+8.4};
        \node[cell, fill=white] at (3.6, -1.2) {-2.1};
        \node[cell, fill=white] at (4.4, -1.2) {-0.9};

        % Row 2: blue
        \node[qlabel] at (-0.1, -2.0) {blue};
        \node[cell, fill=white] at (0.4, -2.0) {-1.0};
        \node[cell, fill=white] at (1.2, -2.0) {+1.9};
        \node[cell, fill=white] at (2.0, -2.0) {+3.5};
        \node[cell, fill=white] at (2.8, -2.0) {+9.1};
        \node[cell, fill=white] at (3.6, -2.0) {-1.4};
        \node[cell, fill=white] at (4.4, -2.0) {-0.3};

        % Row 3: creature (Highlighted Query)
        \node[qlabel, text=politoNavy] at (-0.1, -2.8) {\textbf{creature}};
        \draw[-Stealth, thick, color=politoNavy] (-1.1, -2.8) -- (-0.8, -2.8);
        \node[cell, fill=white] at (0.4, -2.8) {-5.2};
        \node[hcell, text=examAmber] at (1.2, -2.8) {\textbf{+93.0}};
        \node[hcell, text=examAmber] at (2.0, -2.8) {\textbf{+93.4}};
        \node[cell, fill=white] at (2.8, -2.8) {+4.1};
        \node[cell, fill=white] at (3.6, -2.8) {-12.8};
        \node[cell, fill=white] at (4.4, -2.8) {-2.5};

        % Row 4: roamed
        \node[qlabel] at (-0.1, -3.6) {roamed};
        \node[cell, fill=white] at (0.4, -3.6) {-0.4};
        \node[cell, fill=white] at (1.2, -3.6) {-1.8};
        \node[cell, fill=white] at (2.0, -3.6) {-1.2};
        \node[cell, fill=white] at (2.8, -3.6) {+15.2};
        \node[cell, fill=white] at (3.6, -3.6) {+5.0};
        \node[cell, fill=white] at (4.4, -3.6) {+12.6};

        % Row 5: forest
        \node[qlabel] at (-0.1, -4.4) {forest};
        \node[cell, fill=white] at (0.4, -4.4) {-2.1};
        \node[cell, fill=white] at (1.2, -4.4) {-0.7};
        \node[cell, fill=white] at (2.0, -4.4) {-0.3};
        \node[cell, fill=white] at (2.8, -4.4) {+3.2};
        \node[cell, fill=white] at (3.6, -4.4) {+8.9};
        \node[cell, fill=white] at (4.4, -4.4) {+6.4};

        % Highlight bounding box around the creature row matches
        \draw[examAmber, very thick, rounded corners=2pt] (0.8, -3.2) rectangle (2.4, -2.4);
        \node[font=\tiny\bfseries, text=examAmber, anchor=north] at (1.6, -3.3) {Modifiers: $+93.0$ and $+93.4$};
    \end{tikzpicture}%
    }
    \caption{Query-Key compatibility matrix ($\mQ \mK^\top$)}
    \label{fig:3b1b_attn_scores}
\end{subfigure}\hfill
\begin{subfigure}[b]{0.48\textwidth}
    \centering
    \resizebox{\textwidth}{!}{%
    \begin{tikzpicture}[
        >=Stealth,
        font=\small\sffamily,
        vbox/.style={rectangle, draw=politoNavy, fill=skyBlue!30, thick, rounded corners=3pt, minimum width=2.4cm, minimum height=0.7cm, font=\scriptsize\sffamily, align=center},
        wbox/.style={circle, draw=deepdiveTealLine, fill=deepdiveTealBg, thick, minimum size=0.8cm, font=\scriptsize\bfseries\sffamily},
        sumbox/.style={circle, draw=politoNavy, fill=white, thick, minimum size=0.7cm, font=\bfseries},
        dispbox/.style={rectangle, draw=examAmberLine, fill=examAmberBg, thick, rounded corners=4pt, minimum width=2.6cm, minimum height=0.8cm, font=\scriptsize\bfseries\sffamily, align=center}
    ]
        % Softmax weights
        \node[wbox, text=deepdiveTeal] (a1) at (0, 3.2) {$\alpha_2=0.42$};
        \node[vbox] (v1) at (2.8, 3.2) {$\vv_{\text{fluffy}}$ \textmd{\scriptsize(tactile descriptor)}};
        \draw[->, thick, draw=deepdiveTeal] (a1) -- node[above, font=\tiny]{weight} (v1.west);

        \node[wbox, text=deepdiveTeal] (a2) at (0, 1.8) {$\alpha_3=0.58$};
        \node[vbox] (v2) at (2.8, 1.8) {$\vv_{\text{blue}}$ \textmd{\scriptsize(color descriptor)}};
        \draw[->, thick, draw=deepdiveTeal] (a2) -- node[above, font=\tiny]{weight} (v2.west);

        % Weighted sum node
        \node[sumbox] (sum) at (2.8, 0.4) {$\Sigma$};
        \draw[->, thick] (v1) -- (sum);
        \draw[->, thick] (v2) -- (sum);

        % Displacement vector
        \node[dispbox, text=examAmber] (delta) at (2.8, -0.8) {$\Delta \ve_{\text{creature}} = 0.42 \vv_{\text{fluffy}} + 0.58 \vv_{\text{blue}}$};
        \draw[->, very thick, examAmber] (sum) -- (delta);

        % Residual addition
        \node[sumbox] (add) at (2.8, -1.9) {$\oplus$};
        \node[font=\scriptsize\sffamily, text=bodyDark, anchor=east] (orig) at (1.0, -1.9) {$\ve_{\text{creature}}$ (static)};
        \draw[->, thick] (orig) -- (add);
        \draw[->, very thick, examAmber] (delta) -- (add);

        % Output
        \node[rectangle, draw=thmGreenLine, fill=thmGreenBg, thick, rounded corners=3pt, minimum width=2.8cm, minimum height=0.7cm, font=\scriptsize\bfseries\sffamily, text=thmGreen, anchor=north] (out) at (2.8, -2.5) {$\ve_{\text{creature}}'$ (contextual representation)};
        \draw[->, very thick, thmGreen] (add) -- (out);

    \end{tikzpicture}%
    }
    \caption{Value matrix weighting and vector update ($\Delta \mathbf{e}$)}
    \label{fig:3b1b_value_table}
\end{subfigure}
\caption{Geometric attention dynamics and matrix transformations in 3Blue1Brown's pedagogical breakdown \cite{sanderson2024attention}. (a) Pairwise Query-Key dot products evaluate compatibility across all token pairs in \emph{``A fluffy blue creature roamed the verdant forest''}; the Query for \emph{``creature''} strongly matches the Keys for descriptive modifiers \emph{``fluffy''} ($+93.0$) and \emph{``blue''} ($+93.4$). (b) The attention table routes the corresponding Value vectors: after softmax normalization ($0.42$ and $0.58$), the weighted sum produces a displacement vector $\Delta \mathbf{E}_4 = 0.42 \mathbf{V}_2 + 0.58 \mathbf{V}_3$. Adding $\Delta \mathbf{E}_4$ to the generic creature vector $\mathbf{E}_4$ displaces the representation into the specialized manifold coordinates of a fluffy blue creature.}
\label{fig:3b1b_query_key_creature}
\end{figure}

\clearpage
\section{Scaled Dot-Product Attention}
\label{sec:scaled_dot_product_attention}

\subsection{Definition of the scaled attention mechanism}
\label{subsec:scaled_dot_product_defn}

When computing alignment scores between queries and keys, standard dot product $s = \vq^\top \vk$ suffers from a critical numerical failure mode as vector dimensionality increases. To resolve this, the landmark formulation in \cite{vaswani2017} introduced \textbf{Scaled Dot-Product Attention}.

\begin{definition}{Scaled dot-product attention}{scaled_dot_product_defn}
Let $\mQ \in \R^{L_Q \times d_k}$, $\mK \in \R^{L_K \times d_k}$, and $\mV \in \R^{L_K \times d_v}$ denote Query, Key, and Value matrices. The Scaled Dot-Product Attention operator evaluates:
\begin{equation}
\label{eq:scaled_dot_product_matrix}
\text{\texttt{Attention}}(\mQ, \mK, \mV) = \softmax\left( \frac{\mQ \mK^\top}{\sqrt{d_k}} + \mM \right) \mV,
\end{equation}
where $\mM \in \{0, -\infty\}^{L_Q \times L_K}$ is an optional masking matrix, and $\softmax$ is applied row-wise.
\end{definition}

\begin{figure}[htbp]
\centering
\begin{tikzpicture}[
    scale=0.9, every node/.style={transform shape},
    block/.style={draw=politoNavy, fill=skyBlue!30, thick, rounded corners=3pt, minimum width=2.6cm, minimum height=0.7cm, align=center, font=\small\sffamily},
    matblock/.style={draw=examAmberLine, fill=examAmberBg, thick, rounded corners=3pt, minimum width=2.6cm, minimum height=0.7cm, align=center, font=\small\sffamily\bfseries},
    arrow/.style={-{Stealth[length=5pt]}, thick, draw=politoNavy}
]
    \node[font=\sffamily\bfseries] (q_in) at (-2, 0) {Query $\mQ$};
    \node[font=\sffamily\bfseries] (k_in) at (0, 0) {Key $\mK$};
    \node[font=\sffamily\bfseries] (v_in) at (2.5, 0) {Value $\mV$};

    \node[matblock] (matmul1) at (-1, 1.3) {MatMul};
    \draw[arrow] (q_in) |- (matmul1.west);
    \draw[arrow] (k_in) |- (matmul1.east);

    \node[block] (scale) at (-1, 2.4) {Scale ($1/\sqrt{d_k}$)};
    \draw[arrow] (matmul1) -- (scale);

    \node[block, fill=examAmberBg!50] (mask) at (-1, 3.5) {Mask (Optional)};
    \draw[arrow] (scale) -- (mask);

    \node[block] (smax) at (-1, 4.6) {Softmax};
    \draw[arrow] (mask) -- (smax);

    \node[matblock] (matmul2) at (0.8, 5.8) {MatMul};
    \draw[arrow] (smax) |- (matmul2.west);
    \draw[arrow] (v_in) -- ++(0, 5.8) -- (matmul2.east);

    \node[above=0.8cm of matmul2, font=\sffamily\bfseries\color{politoNavy}] (out) {Output $\mO \in \R^{L_Q \times d_v}$};
    \draw[arrow] (matmul2) -- (out);
\end{tikzpicture}
\caption{Computational dataflow of Scaled Dot-Product Attention \cite{vaswani2017}. The dot-product matrix is scaled by $1/\sqrt{d_k}$ to prevent gradient saturation before row-wise softmax normalization and value multiplication.}
\label{fig:scaled_dot_product_tikz}
\end{figure}

\subsection{Why linear projections are mandatory}
\label{subsec:why_linear_projections}

In self-attention, input tokens start as embeddings $\vx_i \in \R^{\dmodel}$. If we computed dot products directly between raw token embeddings $\vx_i^\top \vx_j$, two critical limitations would cripple model expressivity:
\begin{enumerate}
    \item \textbf{Symmetry trap}: Because the dot product is commutative ($\vx_i^\top \vx_j = \vx_j^\top \vx_i$), the unprojected affinity between token $i$ and token $j$ would be strictly symmetric. A token could not express an asymmetric relationship, such as a verb actively seeking its direct object without the direct object reciprocating with identical strength.
    \item \textbf{Role collapse}: A token vector would be forced to play three fundamentally distinct roles simultaneously: formulating a search request (Query), advertising its semantic contents (Key), and providing the contextual payload (Value).
\end{enumerate}
To break this symmetry and enable functional specialization, the model projects each input vector into three distinct latent subspaces using learned weight matrices:
\begin{equation}
\label{eq:qkv_linear_projections}
\mQ = \mX \mW_Q, \quad \mK = \mX \mW_K, \quad \mV = \mX \mW_V,
\end{equation}
where $\mW_Q \in \R^{\dmodel \times d_k}$, $\mW_K \in \R^{\dmodel \times d_k}$, and $\mW_V \in \R^{\dmodel \times d_v}$.

\subsection{Complete numerical walkthrough: reproducing the lecture calculations}
\label{subsec:complete_numerical_walkthrough}

To illustrate these matrix projections and attention operations concretely, we walk through step-by-step the complete numerical worked example, tracking intermediate values down to four decimal places.


Consider an input sequence of $L = 3$ tokens in a 2-dimensional latent embedding space ($\dmodel = 2$), with projection dimensions $d_k = d_v = 2$:
\begin{equation}
\mX = \begin{bmatrix}
0.1 & 0.5 \\
0.1 & 0.2 \\
0.9 & 0.9
\end{bmatrix} \in \R^{3 \times 2}, \quad
\mW_Q = \begin{bmatrix}
0.9 & 0.9 \\
0.5 & 0.3
\end{bmatrix}, \quad
\mW_K = \begin{bmatrix}
1.0 & 0.9 \\
1.0 & 0.6
\end{bmatrix}, \quad
\mW_V = \begin{bmatrix}
0.8 & 0.9 \\
0.9 & 1.0
\end{bmatrix}.
\end{equation}

\subsubsection{Linear projection into Queries, Keys, and Values}
Multiplying $\mX$ by each projection matrix yields:
\begin{equation}
\mQ = \mX \mW_Q = \begin{bmatrix}
0.1(0.9) + 0.5(0.5) & 0.1(0.9) + 0.5(0.3) \\
0.1(0.9) + 0.2(0.5) & 0.1(0.9) + 0.2(0.3) \\
0.9(0.9) + 0.9(0.5) & 0.9(0.9) + 0.9(0.3)
\end{bmatrix} = \begin{bmatrix}
0.34 & 0.24 \\
0.19 & 0.15 \\
1.26 & 1.08
\end{bmatrix},
\end{equation}
\begin{equation}
\mK = \mX \mW_K = \begin{bmatrix}
0.1(1.0) + 0.5(1.0) & 0.1(0.9) + 0.5(0.6) \\
0.1(1.0) + 0.2(1.0) & 0.1(0.9) + 0.2(0.6) \\
0.9(1.0) + 0.9(1.0) & 0.9(0.9) + 0.9(0.6)
\end{bmatrix} = \begin{bmatrix}
0.60 & 0.39 \\
0.30 & 0.21 \\
1.80 & 1.35
\end{bmatrix},
\end{equation}
\begin{equation}
\mV = \mX \mW_V = \begin{bmatrix}
0.1(0.8) + 0.5(0.9) & 0.1(0.9) + 0.5(1.0) \\
0.1(0.8) + 0.2(0.9) & 0.1(0.9) + 0.2(1.0) \\
0.9(0.8) + 0.9(0.9) & 0.9(0.9) + 0.9(1.0)
\end{bmatrix} = \begin{bmatrix}
0.53 & 0.59 \\
0.26 & 0.29 \\
1.53 & 1.71
\end{bmatrix}.
\end{equation}

\subsubsection{Scaled attention logits computation}
Evaluating the full matrix multiplication $\mQ \mK^\top \in \R^{3 \times 3}$ directly:
\begin{equation}
\mQ \mK^\top = \begin{bmatrix}
0.34 & 0.24 \\
0.19 & 0.15 \\
1.26 & 1.08
\end{bmatrix}
\begin{bmatrix}
0.60 & 0.30 & 1.80 \\
0.39 & 0.21 & 1.35
\end{bmatrix}
= \begin{bmatrix}
0.2976 & 0.1524 & 0.9360 \\
0.1725 & 0.0885 & 0.5445 \\
1.1772 & 0.6048 & 3.7260
\end{bmatrix}.
\end{equation}
Alongside this matrix computation, examining row 1 reveals the individual scalar dot-product compatibility scores for the first token ($\vq_1 = [0.34, 0.24]$):
\begin{align}
s_{1,1} &= \vq_1^\top \vk_1 = 0.34(0.60) + 0.24(0.39) = 0.2040 + 0.0936 = 0.2976, \\
s_{1,2} &= \vq_1^\top \vk_2 = 0.34(0.30) + 0.24(0.21) = 0.1020 + 0.0504 = 0.1524, \\
s_{1,3} &= \vq_1^\top \vk_3 = 0.34(1.80) + 0.24(1.35) = 0.6120 + 0.3240 = 0.9360.
\end{align}
Next, we scale the complete matrix by $\frac{1}{\sqrt{d_k}} = \frac{1}{\sqrt{2}} \approx 0.7071$:
\begin{equation}
\label{eq:matrix_scaled_scores_numerical}
\mS = \frac{\mQ \mK^\top}{\sqrt{2}} = \begin{bmatrix}
\frac{0.2976}{\sqrt{2}} & \frac{0.1524}{\sqrt{2}} & \frac{0.9360}{\sqrt{2}} \\[4pt]
\frac{0.1725}{\sqrt{2}} & \frac{0.0885}{\sqrt{2}} & \frac{0.5445}{\sqrt{2}} \\[4pt]
\frac{1.1772}{\sqrt{2}} & \frac{0.6048}{\sqrt{2}} & \frac{3.7260}{\sqrt{2}}
\end{bmatrix}
\approx \begin{bmatrix}
0.2104 & 0.1078 & 0.6619 \\
0.1220 & 0.0626 & 0.3850 \\
0.8324 & 0.4277 & 2.6347
\end{bmatrix}.
\end{equation}

\subsubsection{Matrix softmax normalization onto the probability simplex}
Applying the row-wise softmax operator $\mA = \softmax(\mS)$ exponentiates each row and normalizes by its partition sum.
For row 1:
\begin{align}
\exp(0.2104) &\approx 1.2342, \quad \exp(0.1078) \approx 1.1138, \quad \exp(0.6619) \approx 1.9385, \\
\sum_{j=1}^3 \exp(\tilde{s}_{1,j}) &= 1.2342 + 1.1138 + 1.9385 = 4.2865, \\
\alpha_{1,1} &= \frac{1.2342}{4.2865} \approx 0.2879, \quad \alpha_{1,2} = \frac{1.1138}{4.2865} \approx 0.2598, \quad \alpha_{1,3} = \frac{1.9385}{4.2865} \approx 0.4522.
\end{align}
Evaluating rows 2 and 3 identically yields the full row-stochastic attention weight matrix:
\begin{equation}
\label{eq:matrix_attention_weights_numerical}
\mA = \softmax\left(\frac{\mQ \mK^\top}{\sqrt{d_k}}\right) = \begin{bmatrix}
0.2879 & 0.2598 & 0.4522 \\
0.3083 & 0.2906 & 0.4011 \\
0.1294 & 0.0863 & 0.7843
\end{bmatrix},
\end{equation}
where each row sums precisely to $1.0000$.

\subsubsection{Matrix value aggregation}
The full sequence output matrix is obtained by matrix multiplication $\mO = \mA \mV \in \R^{3 \times 2}$:
\begin{equation}
\label{eq:full_output_matrix_lecture}
\mO = \mA \mV = \begin{bmatrix}
0.2879 & 0.2598 & 0.4522 \\
0.3083 & 0.2906 & 0.4011 \\
0.1294 & 0.0863 & 0.7843
\end{bmatrix}
\begin{bmatrix}
0.53 & 0.59 \\
0.26 & 0.29 \\
1.53 & 1.71
\end{bmatrix}
\approx \begin{bmatrix}
0.91 & 1.02 \\
0.85 & 0.95 \\
1.29 & 1.44
\end{bmatrix}.
\end{equation}
For token 1, this matrix product directly computes the linear combination of value rows:
\begin{equation}
\vo_1 = 0.2879 \begin{bmatrix} 0.53 \\ 0.59 \end{bmatrix} + 0.2598 \begin{bmatrix} 0.26 \\ 0.29 \end{bmatrix} + 0.4522 \begin{bmatrix} 1.53 \\ 1.71 \end{bmatrix} \approx \begin{bmatrix} 0.91 \\ 1.02 \end{bmatrix}.
\end{equation}
This reproduces identically the exact tensor results displayed in the calculations.

\subsection{Dimensional constraints across projection matrices}
\label{subsec:dimensional_constraints}

The dimensional compatibility of the projection matrices is governed by three fundamental algebraic constraints:
\begin{itemize}
    \item \textbf{Queries and Keys must share identical dimensionality ($d_k$)}: The dot product $\vq^\top \vk$ is an inner product between two vectors. An inner product is mathematically undefined unless both vectors live in the identical Euclidean vector space $\R^{d_k}$. Therefore, $\mW_Q$ and $\mW_K$ must always share identical column dimension $d_k$.
    \item \textbf{Values can possess independent dimensionality ($d_v$)}: Value vectors are never multiplied with queries or keys. They are only scaled by scalar attention weights $\alpha_j \in [0, 1]$ and summed together. Consequently, $d_v$ could mathematically differ from $d_k$. In standard architectures (including \cite{vaswani2017}), practitioners choose $d_k = d_v = \dmodel / h$ purely for architectural symmetry and implementation simplicity.
    \item \textbf{Sequence lengths in self-attention vs. cross-attention}: In self-attention, $L_Q = L_K = L_V = L$. In cross-attention, queries stem from the decoder ($L_Q = L_{\text{dec}}$), while keys and values stem from the encoder ($L_K = L_V = L_{\text{enc}}$). The dot-product matrix $\mQ \mK^\top$ has shape $L_{\text{dec}} \times L_{\text{enc}}$, and multiplying by $\mV \in \R^{L_{\text{enc}} \times d_v}$ yields the required decoder output shape $L_{\text{dec}} \times d_v$.
\end{itemize}

\subsection{Mathematical derivation: theoretical necessity of the \texorpdfstring{$\frac{1}{\sqrt{d_k}}$}{1/sqrt(dk)} scaling factor}
\label{subsec:scaling_factor_derivation}

Why must the dot product $\vq^\top \vk$ be scaled specifically by $\frac{1}{\sqrt{d_k}}$? To understand the mathematical necessity of this normalization, consider the variance behavior of unscaled dot products as dimensionality expands.


Assume that the components of query vector $\vq \in \R^{d_k}$ and key vector $\vk \in \R^{d_k}$ are independent, identically distributed (i.i.d.) random variables with zero mean and unit variance:
\begin{equation}
\E[q_i] = 0, \quad \Var(q_i) = 1, \quad \E[k_i] = 0, \quad \Var(k_i) = 1, \quad \forall i \in \{1, 2, \dots, d_k\}.
\end{equation}
The unscaled dot product is the random variable $z = \vq^\top \vk = \sum_{i=1}^{d_k} q_i k_i$.
By linearity of expectation and independence:
\begin{equation}
\E[z] = \E\left[\sum_{i=1}^{d_k} q_i k_i\right] = \sum_{i=1}^{d_k} \E[q_i]\E[k_i] = \sum_{i=1}^{d_k} (0 \cdot 0) = 0.
\end{equation}
To compute the variance of $z$, we expand the definition $\Var(z) = \sum_{i=1}^{d_k} \Var(q_i k_i)$. Using identity $\Var(X Y) = \E[X^2]\E[Y^2] - (\E[X]\E[Y])^2$:
\begin{align}
\Var(q_i k_i) &= \E[q_i^2]\E[k_i^2] - (\E[q_i]\E[k_i])^2 \nonumber \\
&= \left(\Var(q_i) + (\E[q_i])^2\right)\left(\Var(k_i) + (\E[k_i])^2\right) - 0 \nonumber \\
&= (1 + 0)(1 + 0) = 1.
\end{align}
Summing over all $d_k$ independent dimensions:
\begin{equation}
\label{eq:variance_unscaled_dot_product}
\Var(z) = \sum_{i=1}^{d_k} \Var(q_i k_i) = \sum_{i=1}^{d_k} 1 = d_k.
\end{equation}
Thus, the standard deviation of the dot product is $\sigma_z = \sqrt{d_k}$.


For large dimensions (e.g., $d_k = 64$ or $128$), the variance $d_k$ is substantial, causing dot products to take values with magnitude $|z| \gg 1$.
Recall that the derivative of the softmax function with respect to logit $z_i$ is given by:
\begin{equation}
\label{eq:softmax_derivative}
\frac{\partial \softmax(z)_i}{\partial z_j} = \softmax(z)_i \left(\delta_{ij} - \softmax(z)_j\right).
\end{equation}
When logits $z$ have large absolute variance, the exponential $\exp(z_m)$ of the maximal logit completely dominates the sum. The softmax output degenerates into a nearly one-hot indicator distribution ($\alpha_m \approx 1$, and $\alpha_{j \ne m} \approx 0$). In this regime:
\begin{equation}
\frac{\partial \softmax(z)_i}{\partial z_j} \approx 0, \quad \forall i, j.
\end{equation}
The gradients vanish, arresting backpropagation through the attention module.
By scaling $z$ by scalar factor $\frac{1}{\sqrt{d_k}}$, we renormalize the variance back to unity:
\begin{equation}
\label{eq:variance_scaled_dot_product}
\Var\left(\frac{z}{\sqrt{d_k}}\right) = \frac{1}{(\sqrt{d_k})^2}\Var(z) = \frac{1}{d_k} \cdot d_k = 1.
\end{equation}
Renormalizing the variance preserves healthy gradient dynamics across arbitrarily deep networks.

\clearpage
\section{Attention Mechanism Types}
\label{sec:triad_of_attention}

Within the complete Transformer architecture, attention is deployed across three distinct operational settings \cite{vaswani2017}, summarized in Figure~\ref{fig:attention_triad_types}.

\begin{figure}[htbp]
\centering
\resizebox{0.96\textwidth}{!}{%
\begin{tikzpicture}[
    >=Stealth,
    font=\small\sffamily,
    panelbox/.style={draw=lightGrayLine, rounded corners=6pt, fill=white, inner sep=8pt},
    header/.style={font=\bfseries\small, align=center},
    cell/.style={rectangle, draw=white, minimum size=7.5mm, thick, font=\scriptsize\bfseries},
    encact/.style={cell, fill=steelBlue!35, text=politoNavy},
    decact/.style={cell, fill=examAmberLine!40, text=examAmber},
    decmask/.style={cell, fill=bodyDark!15, text=bodyGray!70},
    crossact/.style={cell, fill=deepdiveTealLine!35, text=deepdiveTeal},
    lbl/.style={font=\scriptsize, align=center, text=bodyDark}
]

% Panel 1: Encoder Self-Attention
\node[panelbox, minimum width=4.8cm, minimum height=6.2cm] (p1) at (0,0) {};
\node[header, text=politoNavy] at (0, 2.4) {Encoder Self-Attention\\\textmd{\footnotesize (Bidirectional)}};
\foreach \i in {1,...,4} {
    \foreach \j in {1,...,4} {
        \node[encact] at (-1.5 + \j*0.6, 1.7 - \i*0.6) {$\alpha_{\i\j}$};
    }
}
\node[lbl] at (0, -1.2) {$\mQ \leftarrow \mX \mW_Q, \;\; \mK \leftarrow \mX \mW_K, \;\; \mV \leftarrow \mX \mW_V$};
\node[lbl, text=steelBlue] at (0, -1.8) {\textbf{Unconstrained context}:\\All tokens attend to all positions};

% Panel 2: Decoder Masked Self-Attention
\node[panelbox, minimum width=4.8cm, minimum height=6.2cm] (p2) at (5.4,0) {};
\node[header, text=examAmber] at (5.4, 2.4) {Decoder Masked Attention\\\textmd{\footnotesize (Causal Autoregressive)}};
\foreach \i in {1,...,4} {
    \foreach \j in {1,...,4} {
        \pgfmathparse{\j<=\i?1:0}
        \ifnum\pgfmathresult=1
            \node[decact] at (3.9 + \j*0.6, 1.7 - \i*0.6) {$\alpha_{\i\j}$};
        \else
            \node[decmask] at (3.9 + \j*0.6, 1.7 - \i*0.6) {$-\infty$};
        \fi
    }
}
\node[lbl] at (5.4, -1.2) {Causal mask: $\mM_{ij} = -\infty$ for $j > i$};
\node[lbl, text=examAmber] at (5.4, -1.8) {\textbf{Causal constraint}:\\Tokens only attend to past and present};

% Panel 3: Cross-Attention
\node[panelbox, minimum width=4.8cm, minimum height=6.2cm] (p3) at (10.8,0) {};
\node[header, text=deepdiveTeal] at (10.8, 2.4) {Encoder-Decoder Cross-Attention\\\textmd{\footnotesize (Conditioning \& Grounding)}};
\foreach \i in {1,...,3} {
    \foreach \j in {1,...,4} {
        \node[crossact] at (9.3 + \j*0.6, 1.5 - \i*0.6) {$\alpha_{\i\j}$};
    }
}
\node[lbl] at (10.8, -1.2) {$\mQ \leftarrow \text{Decoder}, \quad (\mK, \mV) \leftarrow \text{Encoder}$};
\node[lbl, text=deepdiveTeal] at (10.8, -1.8) {\textbf{Memory bridge}:\\Decoder queries attend to full source sequence};

\end{tikzpicture}%
}
\caption{The triad of attention mechanisms across the Transformer architecture. (a) Encoder self-attention processes inputs bidirectionally with an unconstrained attention matrix. (b) Decoder masked self-attention applies a lower-triangular causal mask ($\mM_{ij} = -\infty$ for $j > i$) to enforce autoregressive generation without future information leakage. (c) Encoder-decoder cross-attention allows decoder queries to attend across all encoder key-value pairs, bridging representations across different sequence lengths.}
\label{fig:attention_triad_types}
\end{figure}

\subsection{Encoder self-attention: unconstrained bidirectional contextualization}
\label{subsec:encoder_self_attention}

\begin{figure}[htbp]
\centering
\begin{tikzpicture}[
    >=Stealth,
    font=\small\sffamily,
    block/.style={rectangle, draw=tfOrangeBorder, fill=tfOrange, thick, rounded corners=4pt, minimum width=3.8cm, minimum height=0.9cm, align=center, font=\small\bfseries},
    normblock/.style={rectangle, draw=tfYellowBorder, fill=tfYellow, thick, rounded corners=4pt, minimum width=3.8cm, minimum height=0.7cm, align=center, font=\small\bfseries},
    projnode/.style={rectangle, draw=steelBlue, fill=skyBlue!40, thick, rounded corners=3pt, minimum width=0.9cm, minimum height=0.6cm, align=center, font=\scriptsize\bfseries},
    sum/.style={circle, draw=politoNavy, fill=white, thick, minimum size=0.6cm, font=\bfseries},
    box/.style={draw=lightGrayLine, fill=tfGrayBg, rounded corners=6pt, inner sep=12pt}
]
    % Outer container
    \node[box, minimum width=6.2cm, minimum height=7.2cm] (frame) at (0, 3.2) {};
    \node[font=\footnotesize\bfseries\scshape, text=politoNavy, anchor=north west] at (frame.north west) [xshift=6pt, yshift=-6pt] {Encoder Layer: Self-Attention Sublayer};

    % Input
    \node (in) at (0, 0) {Input $\mX \in \R^{L \times \dmodel}$};

    % Split node
    \coordinate (split) at (0, 0.7);
    \draw[thick] (in) -- (split);

    % Projections Q, K, V
    \node[projnode] (qproj) at (-1.4, 1.4) {$\mW_Q$};
    \node[projnode] (kproj) at (0, 1.4) {$\mW_K$};
    \node[projnode] (vproj) at (1.4, 1.4) {$\mW_V$};

    \draw[->, thick] (split) -| (qproj.south);
    \draw[->, thick] (split) -- (kproj.south);
    \draw[->, thick] (split) -| (vproj.south);

    % MHA Block
    \node[block] (mha) at (0, 2.7) {Multi-Head Attention\\[-2pt]\textmd{\scriptsize 8 parallel heads $\cdot$ Unconstrained ($\mM=\mathbf{0}$)}};
    \draw[->, thick] (qproj) -- node[left, font=\scriptsize]{$\mQ$} (-1.4, 2.25);
    \draw[->, thick] (kproj) -- node[right, font=\scriptsize]{$\mK$} (0, 2.25);
    \draw[->, thick] (vproj) -- node[right, font=\scriptsize]{$\mV$} (1.4, 2.25);

    % Linear output projection
    \node[projnode, minimum width=2.4cm] (wout) at (0, 3.8) {Linear Projection $\mW_O$};
    \draw[->, thick] (mha) -- (wout);

    % Add node
    \node[sum] (add) at (0, 4.8) {$\oplus$};
    \draw[->, thick] (wout) -- (add);

    % Residual shortcut
    \draw[->, thick, rounded corners=5pt, draw=crimsonRed] (split) -- ++(-2.4, 0) -- ++(0, 4.1) -- node[above, font=\tiny\bfseries, text=crimsonRed]{Residual Shortcut} (add.west);

    % LayerNorm
    \node[normblock] (ln) at (0, 5.8) {Add \& Norm (LayerNorm)};
    \draw[->, thick] (add) -- (ln);

    % Output
    \node (out) at (0, 7.0) {Output to FFN: $\mH^{(1)} \in \R^{L \times \dmodel}$};
    \draw[->, very thick] (ln) -- (out);

\end{tikzpicture}
\caption{Architectural schematic of the Encoder Self-Attention sublayer \cite{vaswani2017}. Input tokens $\mathbf{X}$ are projected into queries $\mathbf{Q}$, keys $\mathbf{K}$, and values $\mathbf{V}$. Unconstrained multi-head attention evaluates all-to-all token visibility without masking ($\mathbf{M}=\mathbf{0}$), followed by linear output projection $\mathbf{W}_O$, identity residual addition, and Layer Normalization.}
\label{fig:vaswani_highlight_encoder_mha}
\end{figure}


In natural language understanding, the meaning of a word is profoundly shaped by tokens that appear both \emph{before} and \emph{after} it. For example, consider syntactic ambiguity, pronoun reference, or polysemy: in the sentence \emph{``The bank of the river overflowed''}, the polysemous noun \emph{``bank''} cannot be resolved into a geological landform rather than a financial institution without inspecting the downstream word \emph{``river''}. 
Because the Encoder's sole objective is representation learning rather than autoregressive word generation, there is no risk of label leakage. Therefore, the Encoder must have access to the entire sentence simultaneously: the attention matrix $\mA \in \R^{L_{\text{enc}} \times L_{\text{enc}}}$ is completely unconstrained ($\mM = \mat{0}$), granting every token an instantaneous $\mathcal{O}(1)$ global receptive field across all $L_{\text{enc}}$ positions.

In the Encoder, queries, keys, and values are all linearly projected from the preceding encoder layer's representations:
\begin{equation}
\label{eq:encoder_self_attn_projections}
\mQ = \mX \mW_Q, \quad \mK = \mX \mW_K, \quad \mV = \mX \mW_V,
\end{equation}
where $\mX \in \R^{L_{\text{enc}} \times \dmodel}$, $\mW_Q, \mW_K \in \R^{\dmodel \times d_k}$, and $\mW_V \in \R^{\dmodel \times d_v}$. No masking matrix is applied ($\mM = \mat{0}$). Every token $i$ can attend freely to all other tokens $j \in \{1, 2, \dots, L_{\text{enc}}\}$, producing unconstrained bidirectional representations.

\subsubsection{Worked example: bidirectional polysemy disambiguation}
Consider an English input sentence of $L_{\text{enc}} = 3$ tokens:
\begin{equation}
X = [\text{``bank''}_1,\, \text{``of''}_2,\, \text{``river''}_3].
\end{equation}
The token \emph{``bank''} is polysemous. Without inspecting the future token \emph{``river''}, disambiguation is impossible.
Let $d_k = d_v = 2$. Suppose the linear projections $\mX \mW_Q$, $\mX \mW_K$, and $\mX \mW_V$ yield:
\begin{equation}
\label{eq:enc_example_matrices}
\mQ = \begin{bmatrix}
1.0 & 0.5 \\
0.2 & 0.8 \\
0.9 & 1.2
\end{bmatrix}, \quad
\mK = \begin{bmatrix}
0.8 & 0.4 \\
0.1 & 0.9 \\
1.0 & 1.1
\end{bmatrix}, \quad
\mV = \begin{bmatrix}
1.5 & 0.2 \\
0.1 & 0.8 \\
0.4 & 2.0
\end{bmatrix}.
\end{equation}
We execute the attention mechanism step by step:
\begin{enumerate}
    \item \textbf{Raw dot-product computation}:
    Evaluating unscaled inner products $\mQ \mK^\top \in \R^{3 \times 3}$:
    \begin{equation}
    \mQ \mK^\top = \begin{bmatrix}
    (1.0)(0.8) + (0.5)(0.4) & (1.0)(0.1) + (0.5)(0.9) & (1.0)(1.0) + (0.5)(1.1) \\
    (0.2)(0.8) + (0.8)(0.4) & (0.2)(0.1) + (0.8)(0.9) & (0.2)(1.0) + (0.8)(1.1) \\
    (0.9)(0.8) + (1.2)(0.4) & (0.9)(0.1) + (1.2)(0.9) & (0.9)(1.0) + (1.2)(1.1)
    \end{bmatrix}.
    \end{equation}
    Evaluating each scalar entry yields:
    \begin{equation}
    \mQ \mK^\top = \begin{bmatrix}
    1.00 & 0.55 & 1.55 \\
    0.48 & 0.74 & 1.08 \\
    1.20 & 1.17 & 2.22
    \end{bmatrix}.
    \end{equation}
    \item \textbf{Scaling by $\sqrt{d_k} = \sqrt{2} \approx 1.414$}:
    \begin{equation}
    \label{eq:enc_scaled_scores}
    \mS = \frac{\mQ \mK^\top}{\sqrt{2}} \approx \begin{bmatrix}
    0.707 & 0.389 & 1.096 \\
    0.339 & 0.523 & 0.764 \\
    0.849 & 0.827 & 1.570
    \end{bmatrix}.
    \end{equation}
    \item \textbf{Bidirectional softmax normalization}:
    Since no mask is present ($\mM = \mat{0}$), we compute row-wise $\softmax$ directly across all columns:
    For row 1 (corresponding to \emph{``bank''}):
    \begin{align}
    \exp(0.707) &\approx 2.028, \quad \exp(0.389) \approx 1.475, \quad \exp(1.096) \approx 2.992, \\
    Z_1 &= 2.028 + 1.475 + 2.992 = 6.495, \\
    \alpha_{11} &= \frac{2.028}{6.495} \approx 0.312, \quad
    \alpha_{12} = \frac{1.475}{6.495} \approx 0.227, \quad
    \alpha_{13} = \frac{2.992}{6.495} \approx 0.461.
    \end{align}
    Applying this across all rows yields the dense attention matrix:
    \begin{equation}
    \label{eq:enc_attn_weights}
    \mA_{\text{enc}} = \begin{bmatrix}
    0.312 & 0.227 & 0.461 \\
    0.268 & 0.322 & 0.410 \\
    0.248 & 0.243 & 0.510
    \end{bmatrix}.
    \end{equation}
    \item \textbf{Contextual value aggregation}:
    Multiplying $\mA_{\text{enc}}$ by $\mV \in \R^{3 \times 2}$ yields the updated contextual representations:
    \begin{equation}
    \label{eq:enc_output_matrix}
    \mO_{\text{enc}} = \mA_{\text{enc}} \mV = \begin{bmatrix}
    0.312 \vv_1 + 0.227 \vv_2 + 0.461 \vv_3 \\
    0.268 \vv_1 + 0.322 \vv_2 + 0.410 \vv_3 \\
    0.248 \vv_1 + 0.243 \vv_2 + 0.510 \vv_3
    \end{bmatrix}
    \approx \begin{bmatrix}
    0.675 & 1.165 \\
    0.598 & 1.131 \\
    0.600 & 1.263
    \end{bmatrix}.
    \end{equation}
\end{enumerate}

Notice that token 1 (\emph{``bank''}) allocates its largest attention weight ($\alpha_{13} = 0.461$) to token 3 (\emph{``river''}), looking \emph{ahead} into future context. This unconstrained bidirectional information exchange allows \emph{``bank''} to instantly absorb the geological semantics of \emph{``river''} into its contextual output $\vo_1$.

\subsection{Decoder masked self-attention: causal autoregressive constraint}
\label{subsec:decoder_masked_self_attention}

\begin{figure}[htbp]
\centering
\begin{tikzpicture}[
    >=Stealth,
    font=\small\sffamily,
    block/.style={rectangle, draw=tfOrangeBorder, fill=tfOrange, thick, rounded corners=4pt, minimum width=3.8cm, minimum height=0.9cm, align=center, font=\small\bfseries},
    normblock/.style={rectangle, draw=tfYellowBorder, fill=tfYellow, thick, rounded corners=4pt, minimum width=3.8cm, minimum height=0.7cm, align=center, font=\small\bfseries},
    projnode/.style={rectangle, draw=steelBlue, fill=skyBlue!40, thick, rounded corners=3pt, minimum width=0.9cm, minimum height=0.6cm, align=center, font=\scriptsize\bfseries},
    sum/.style={circle, draw=politoNavy, fill=white, thick, minimum size=0.6cm, font=\bfseries},
    box/.style={draw=lightGrayLine, fill=tfGrayBg, rounded corners=6pt, inner sep=12pt}
]
    % Outer container
    \node[box, minimum width=6.2cm, minimum height=7.2cm] (frame) at (0, 3.2) {};
    \node[font=\footnotesize\bfseries\scshape, text=examAmber, anchor=north west] at (frame.north west) [xshift=6pt, yshift=-6pt] {Decoder Layer: Masked Self-Attention Sublayer};

    % Input
    \node (in) at (0, 0) {Shifted Targets $\mY \in \R^{L_{\text{dec}} \times \dmodel}$};

    % Split node
    \coordinate (split) at (0, 0.7);
    \draw[thick] (in) -- (split);

    % Projections Q, K, V
    \node[projnode] (qproj) at (-1.4, 1.4) {$\mW_Q$};
    \node[projnode] (kproj) at (0, 1.4) {$\mW_K$};
    \node[projnode] (vproj) at (1.4, 1.4) {$\mW_V$};

    \draw[->, thick] (split) -| (qproj.south);
    \draw[->, thick] (split) -- (kproj.south);
    \draw[->, thick] (split) -| (vproj.south);

    % Masked MHA Block
    \node[block] (mha) at (0, 2.7) {Masked Multi-Head Attention\\[-2pt]\textmd{\scriptsize Causal Mask $\mM_{ij} = -\infty$ for $j > i$}};
    \draw[->, thick] (qproj) -- node[left, font=\scriptsize]{$\mQ$} (-1.4, 2.25);
    \draw[->, thick] (kproj) -- node[right, font=\scriptsize]{$\mK$} (0, 2.25);
    \draw[->, thick] (vproj) -- node[right, font=\scriptsize]{$\mV$} (1.4, 2.25);

    % Linear output projection
    \node[projnode, minimum width=2.4cm] (wout) at (0, 3.8) {Linear Projection $\mW_O$};
    \draw[->, thick] (mha) -- (wout);

    % Add node
    \node[sum] (add) at (0, 4.8) {$\oplus$};
    \draw[->, thick] (wout) -- (add);

    % Residual shortcut
    \draw[->, thick, rounded corners=5pt, draw=examAmberLine] (split) -- ++(-2.4, 0) -- ++(0, 4.1) -- node[above, font=\tiny\bfseries, text=examAmber]{Residual Shortcut} (add.west);

    % LayerNorm
    \node[normblock] (ln) at (0, 5.8) {Add \& Norm (LayerNorm)};
    \draw[->, thick] (add) -- (ln);

    % Output
    \node (out) at (0, 7.0) {Output to Cross-Attn: $\mH_{\text{dec, self}}^{(1)}$};
    \draw[->, very thick] (ln) -- (out);

\end{tikzpicture}
\caption{Architectural schematic of the Decoder Masked Self-Attention sublayer \cite{vaswani2017}. Causal masking enforces autoregressive conditioning by blinding all future positions with $-\infty$, ensuring query $i$ attends exclusively to preceding tokens $j \le i$.}
\label{fig:vaswani_highlight_masked_mha}
\end{figure}


Unlike the Encoder, the Decoder is tasked with sequential language generation: producing the next token $y_t$ conditioned strictly on the prefix of tokens emitted so far $[y_1, y_2, \dots, y_{t-1}]$. At test time, future tokens $y_{t+1}, y_{t+2}, \dots$ do not exist.
During training, however, the entire target sequence is supplied to the model simultaneously to enable parallel execution via Teacher Forcing. If unconstrained self-attention were permitted in the decoder, a token at position $t$ could directly attend to position $t+1$ and cheat by copying the ground-truth answer. The model would suffer from substantial information leakage, failing to learn any predictive distribution over future tokens.
To enforce the arrow of time while preserving parallel tensor execution on hardware, a \textbf{causal mask} $\mM$ is added to the attention logits prior to softmax, strictly setting all future attention weights to zero.

To preserve causality while retaining parallel training under Teacher Forcing, a \textbf{causal mask} $\mM \in \R^{L_{\text{dec}} \times L_{\text{dec}}}$ is added to the scaled dot products prior to the softmax operator:
\begin{equation}
\label{eq:causal_mask_definition}
M_{i, j} = \begin{cases}
0 & \text{if } j \le i, \\
-\infty & \text{if } j > i.
\end{cases}
\end{equation}
In matrix form for a sequence of length 3:
\begin{equation}
\mM = \begin{bmatrix}
0 & -\infty & -\infty \\
0 & 0 & -\infty \\
0 & 0 & 0
\end{bmatrix}.
\end{equation}
Because $\exp(-\infty) = 0$, evaluating the row-wise softmax guarantees:
\begin{equation}
\alpha_{i, j} = 0 \quad \forall j > i.
\end{equation}
From a hardware efficiency standpoint, computing the full dense GEMM operation $\mQ \mK^\top$ followed by a tensor addition with $\mM$ fully saturates GPU tensor cores and memory buses, outperforming jagged triangular loops.

\subsubsection{Worked example: autoregressive generation under causal masking}
Now consider the Decoder generating the Italian translation prefix of length $L_{\text{dec}} = 3$:
\begin{equation}
Y = [y_1 = \langle\text{BOS}\rangle,\, y_2 = \text{``riva''},\, y_3 = \text{``del''}].
\end{equation}
All queries, keys, and values are projected from the decoder's target token embeddings ($d_k = d_v = 2$):
\begin{equation}
\label{eq:dec_example_matrices}
\mQ = \begin{bmatrix}
0.5 & 0.5 \\
1.2 & 0.3 \\
0.8 & 1.0
\end{bmatrix}, \quad
\mK = \begin{bmatrix}
0.6 & 0.4 \\
1.0 & 0.2 \\
0.7 & 0.9
\end{bmatrix}, \quad
\mV = \begin{bmatrix}
0.2 & 0.5 \\
1.8 & 0.4 \\
0.6 & 1.4
\end{bmatrix}.
\end{equation}
Let us trace the computation step by step:
\begin{enumerate}
    \item \textbf{Raw scaled dot products (unmasked)}:
    Computing $\mS = \frac{\mQ \mK^\top}{\sqrt{2}}$:
    \begin{equation}
    \label{eq:dec_unmasked_scores}
    \mS = \frac{1}{\sqrt{2}} \begin{bmatrix}
    0.50 & 0.60 & 0.80 \\
    0.84 & 1.26 & 1.11 \\
    0.88 & 1.00 & 1.46
    \end{bmatrix}
    \approx \begin{bmatrix}
    0.354 & 0.424 & 0.566 \\
    0.594 & 0.891 & 0.785 \\
    0.622 & 0.707 & 1.032
    \end{bmatrix}.
    \end{equation}
    Notice that unmasked score $S_{13} = 0.566$ is the highest in the first row: without masking, $\langle\text{BOS}\rangle$ would strongly attend to future token \emph{``del''}.
    \item \textbf{Applying the causal mask $\mM$}:
    We add $\mM$ to $\mS$:
    \begin{equation}
    \label{eq:dec_masked_scores}
    \mS_{\text{masked}} = \mS + \mM = \begin{bmatrix}
    0.354 & -\infty & -\infty \\
    0.594 & 0.891 & -\infty \\
    0.622 & 0.707 & 1.032
    \end{bmatrix}.
    \end{equation}
    \item \textbf{Row-wise softmax under masking}:
    \begin{itemize}
        \item \textbf{Row 1 ($y_1 = \langle\text{BOS}\rangle$)}:
        $\exp(0.354) \approx 1.425$, while $\exp(-\infty) = 0$:
        \begin{equation}
        \alpha_{11} = \frac{1.425}{1.425 + 0 + 0} = 1.000, \quad \alpha_{12} = 0.000, \quad \alpha_{13} = 0.000.
        \end{equation}
        Position 1 can \emph{only} attend to itself.
        \item \textbf{Row 2 ($y_2 = \text{``riva''}$)}:
        $\exp(0.594) \approx 1.811$, $\exp(0.891) \approx 2.438$, $\exp(-\infty) = 0$:
        \begin{align}
        Z_2 &= 1.811 + 2.438 = 4.249, \\
        \alpha_{21} &= \frac{1.811}{4.249} \approx 0.426, \quad
        \alpha_{22} = \frac{2.438}{4.249} \approx 0.574, \quad
        \alpha_{23} = 0.000.
        \end{align}
        Position 2 attends strictly to the prefix $\{y_1, y_2\}$; future token $y_3$ is zeroed out.
        \item \textbf{Row 3 ($y_3 = \text{``del''}$)}:
        $\exp(0.622) \approx 1.863$, $\exp(0.707) \approx 2.028$, $\exp(1.032) \approx 2.807$:
        \begin{equation}
        Z_3 = 1.863 + 2.028 + 2.807 = 6.698 \implies
        \alpha_{31} \approx 0.278, \quad \alpha_{32} \approx 0.303, \quad \alpha_{33} \approx 0.419.
        \end{equation}
    \end{itemize}
    This produces the strictly lower-triangular attention matrix:
    \begin{equation}
    \label{eq:dec_causal_weights}
    \mA_{\text{dec}} = \begin{bmatrix}
    1.000 & 0.000 & 0.000 \\
    0.426 & 0.574 & 0.000 \\
    0.278 & 0.303 & 0.419
    \end{bmatrix}.
    \end{equation}
    \item \textbf{Causal value aggregation}:
    Multiplying by $\mV \in \R^{3 \times 2}$:
    \begin{equation}
    \label{eq:dec_output_matrix}
    \mO_{\text{dec}} = \mA_{\text{dec}} \mV = \begin{bmatrix}
    1.000 \vv_1 \\
    0.426 \vv_1 + 0.574 \vv_2 \\
    0.278 \vv_1 + 0.303 \vv_2 + 0.419 \vv_3
    \end{bmatrix}
    \approx \begin{bmatrix}
    0.200 & 0.500 \\
    1.118 & 0.443 \\
    0.852 & 0.847
    \end{bmatrix}.
    \end{equation}
\end{enumerate}

Contrast this directly with Encoder self-attention: despite $S_{13} = 0.566$ having high compatibility, the causal mask completely extinguishes its influence ($\alpha_{13} = 0$). Output $\vo_1$ is computed strictly from $\vv_1$, preserving the causality required for autoregressive sequence modeling.

\subsection{Encoder-decoder cross-attention: dynamic source grounding}
\label{subsec:cross_attention_mechanics}

\begin{figure}[htbp]
\centering
\begin{tikzpicture}[
    >=Stealth,
    font=\small\sffamily,
    block/.style={rectangle, draw=tfOrangeBorder, fill=tfOrange, thick, rounded corners=4pt, minimum width=3.8cm, minimum height=0.9cm, align=center, font=\small\bfseries},
    normblock/.style={rectangle, draw=tfYellowBorder, fill=tfYellow, thick, rounded corners=4pt, minimum width=3.8cm, minimum height=0.7cm, align=center, font=\small\bfseries},
    projnode/.style={rectangle, draw=steelBlue, fill=skyBlue!40, thick, rounded corners=3pt, minimum width=0.9cm, minimum height=0.6cm, align=center, font=\scriptsize\bfseries},
    encnode/.style={rectangle, draw=deepdiveTealLine, fill=deepdiveTealBg, thick, rounded corners=3pt, minimum width=0.9cm, minimum height=0.6cm, align=center, font=\scriptsize\bfseries, text=deepdiveTeal},
    sum/.style={circle, draw=politoNavy, fill=white, thick, minimum size=0.6cm, font=\bfseries},
    box/.style={draw=lightGrayLine, fill=tfGrayBg, rounded corners=6pt, inner sep=12pt}
]
    % Outer container
    \node[box, minimum width=7.2cm, minimum height=7.2cm] (frame) at (0.4, 3.2) {};
    \node[font=\footnotesize\bfseries\scshape, text=deepdiveTeal, anchor=north west] at (frame.north west) [xshift=6pt, yshift=-6pt] {Decoder Layer: Cross-Attention Sublayer};

    % Inputs
    \node (indec) at (0, 0) {Decoder Input: $\mH_{\text{dec, self}} \in \R^{L_{\text{dec}} \times \dmodel}$};
    \node[anchor=east] (inenc) at (-1.8, 1.8) {Encoder Memory: $\mH_{\text{enc}} \in \R^{L_{\text{enc}} \times \dmodel}$};

    % Decoder query projection
    \coordinate (splitdec) at (0, 0.6);
    \draw[thick] (indec) -- (splitdec);
    \node[projnode] (qproj) at (-0.8, 1.4) {$\mW_Q$};
    \draw[->, thick] (splitdec) -| (qproj.south);

    % Encoder key-value projections
    \coordinate (splitenc) at (-1.5, 1.4);
    \draw[->, thick, draw=deepdiveTeal] (inenc.south east) -- (splitenc);
    \node[encnode] (kproj) at (0.4, 1.4) {$\mW_K$};
    \node[encnode] (vproj) at (1.6, 1.4) {$\mW_V$};
    \draw[->, thick, draw=deepdiveTeal] (splitenc) -| (kproj.south);
    \draw[->, thick, draw=deepdiveTeal] (splitenc) -| (vproj.south);

    % Cross MHA Block
    \node[block] (mha) at (0.4, 2.7) {Cross Multi-Head Attention\\[-2pt]\textmd{\scriptsize $\mQ \leftarrow \text{Decoder}, \quad (\mK, \mV) \leftarrow \text{Encoder}$}};
    \draw[->, thick] (qproj) -- node[left, font=\scriptsize]{$\mQ_{\text{dec}}$} (-0.8, 2.25);
    \draw[->, thick, draw=deepdiveTeal] (kproj) -- node[right, font=\scriptsize]{$\mK_{\text{enc}}$} (0.4, 2.25);
    \draw[->, thick, draw=deepdiveTeal] (vproj) -- node[right, font=\scriptsize]{$\mV_{\text{enc}}$} (1.6, 2.25);

    % Linear output projection
    \node[projnode, minimum width=2.4cm] (wout) at (0.4, 3.8) {Linear Projection $\mW_O$};
    \draw[->, thick] (mha) -- (wout);

    % Add node
    \node[sum] (add) at (0.4, 4.8) {$\oplus$};
    \draw[->, thick] (wout) -- (add);

    % Residual shortcut from decoder stream
    \draw[->, thick, rounded corners=5pt, draw=deepdiveTealLine] (splitdec) -- ++(2.4, 0) -- ++(0, 4.2) -- node[right, font=\tiny\bfseries, text=deepdiveTeal]{Decoder Shortcut} (add.east);

    % LayerNorm
    \node[normblock] (ln) at (0.4, 5.8) {Add \& Norm (LayerNorm)};
    \draw[->, thick] (add) -- (ln);

    % Output
    \node (out) at (0.4, 7.0) {Output to Decoder FFN: $\mH_{\text{dec, cross}} \in \R^{L_{\text{dec}} \times \dmodel}$};
    \draw[->, very thick] (ln) -- (out);

\end{tikzpicture}
\caption{Architectural schematic of the Encoder-Decoder Cross-Attention sublayer \cite{vaswani2017}. The decoder generates Query vectors $\mathbf{Q}$ from its self-attended states, while Key $\mathbf{K}$ and Value $\mathbf{V}$ vectors are projected from the final encoder representations, allowing the decoder to dynamically attend across the entire source sequence.}
\label{fig:vaswani_highlight_cross_mha}
\end{figure}


In sequence-to-sequence translation, the decoder is generating tokens in the target language, while the encoder holds the semantic context of the source sentence.
Recall the fundamental roles of the attention triad:
\begin{itemize}
    \item A \textbf{Query} asks: \emph{``What information am I currently searching for to generate my next word?''} Because the entity generating the next word is the decoder, the queries $\mQ$ must be projected from the decoder's current representation.
    \item A \textbf{Key} advertises: \emph{``What features do I describe in the source sentence?''}
    \item A \textbf{Value} provides: \emph{``What source meaning should be retrieved?''}
\end{itemize}
Because the source content is held exclusively in the encoder, the keys $\mK$ and values $\mV$ must originate from the encoder's output representations. If queries were taken from the encoder instead, the source sentence would be searching the target words, which completely contradicts the directional objective of translating source meaning into target text. By projecting queries from the decoder and keys/values from the encoder, cross-attention acts as a dynamic search engine allowing each generated target token to query the entire uncompressed source sentence.

In cross-attention, the model bridges the target sequence with the source sequence.
The queries originate from the decoder's intermediate representations $\mH_{\text{dec}} \in \R^{L_{\text{dec}} \times \dmodel}$, while keys and values originate from the uncompressed output of the encoder stack $\mH_{\text{enc}} \in \R^{L_{\text{enc}} \times \dmodel}$:
\begin{equation}
\label{eq:cross_attention_projections_explicit}
\mQ = \mH_{\text{dec}} \mW_Q^{\text{cross}} \in \R^{L_{\text{dec}} \times d_k}, \quad
\mK = \mH_{\text{enc}} \mW_K^{\text{cross}} \in \R^{L_{\text{enc}} \times d_k}, \quad
\mV = \mH_{\text{enc}} \mW_V^{\text{cross}} \in \R^{L_{\text{enc}} \times d_v}.
\end{equation}
Crucially:
\begin{itemize}
    \item \textbf{Asymmetric sequence lengths}: The number of queries $L_{\text{dec}}$ does not need to equal the number of keys/values $L_{\text{enc}}$. In translation, sentences in different languages rarely have identical token counts.
    \item \textbf{No causal mask on the source}: The source sentence was fully processed by the encoder before decoding began. Thus, every decoder query position can inspect \emph{all} $L_{\text{enc}}$ source tokens ($\mM = \mat{0}$).
\end{itemize}

\subsubsection{Worked example: cross-lingual alignment with asymmetric lengths}
Suppose our decoder has generated a target prefix of length $L_{\text{dec}} = 2$ ($[\langle\text{BOS}\rangle,\, \text{``riva''}]$), and cross-attends to our earlier English source sentence of length $L_{\text{enc}} = 3$ ($[\text{``bank''},\, \text{``of''},\, \text{``river''}]$).
The decoder provides $L_{\text{dec}} = 2$ queries $\mQ \in \R^{2 \times 2}$ (taken from the first two rows of \eqref{eq:dec_example_matrices}), while the encoder provides $L_{\text{enc}} = 3$ keys and values $\mK, \mV \in \R^{3 \times 2}$ (from \eqref{eq:enc_example_matrices}):
\begin{equation}
\label{eq:cross_example_matrices}
\mQ = \begin{bmatrix}
0.5 & 0.5 \\
1.2 & 0.3
\end{bmatrix}, \quad
\mK = \begin{bmatrix}
0.8 & 0.4 \\
0.1 & 0.9 \\
1.0 & 1.1
\end{bmatrix}, \quad
\mV = \begin{bmatrix}
1.5 & 0.2 \\
0.1 & 0.8 \\
0.4 & 2.0
\end{bmatrix}.
\end{equation}
Let us trace the rectangular cross-attention computation step by step:
\begin{enumerate}
    \item \textbf{Rectangular dot-product matrix ($L_{\text{dec}} \times L_{\text{enc}} = 2 \times 3$)}:
    \begin{equation}
    \mQ \mK^\top = \begin{bmatrix}
    (0.5)(0.8) + (0.5)(0.4) & (0.5)(0.1) + (0.5)(0.9) & (0.5)(1.0) + (0.5)(1.1) \\
    (1.2)(0.8) + (0.3)(0.4) & (1.2)(0.1) + (0.3)(0.9) & (1.2)(1.0) + (0.3)(1.1)
    \end{bmatrix}.
    \end{equation}
    Evaluating the pairwise inner products yields:
    \begin{equation}
    \mQ \mK^\top = \begin{bmatrix}
    0.60 & 0.50 & 1.05 \\
    1.08 & 0.39 & 1.53
    \end{bmatrix}.
    \end{equation}
    \item \textbf{Scaling by $\sqrt{d_k} = \sqrt{2}$}:
    \begin{equation}
    \label{eq:cross_scaled_scores}
    \mS_{\text{cross}} = \frac{\mQ \mK^\top}{\sqrt{2}} \approx \begin{bmatrix}
    0.424 & 0.354 & 0.742 \\
    0.764 & 0.276 & 1.082
    \end{bmatrix} \in \R^{2 \times 3}.
    \end{equation}
    \item \textbf{Row-wise softmax normalization (over $L_{\text{enc}} = 3$ source keys)}:
    \begin{itemize}
        \item \textbf{Row 1 ($\vq_{\langle\text{BOS}\rangle}$)}:
        $\exp(0.424) \approx 1.528$, $\exp(0.354) \approx 1.425$, $\exp(0.742) \approx 2.100$:
        \begin{equation}
        Z_1 = 1.528 + 1.425 + 2.100 = 5.053 \implies
        \mA_{1, :} \approx [0.302, \; 0.282, \; 0.416].
        \end{equation}
        \item \textbf{Row 2 ($\vq_{\text{``riva''}}$)}:
        $\exp(0.764) \approx 2.147$, $\exp(0.276) \approx 1.318$, $\exp(1.082) \approx 2.951$:
        \begin{equation}
        Z_2 = 2.147 + 1.318 + 2.951 = 6.416 \implies
        \mA_{2, :} \approx [0.335, \; 0.205, \; 0.460].
        \end{equation}
    \end{itemize}
    This yields the rectangular cross-attention matrix:
    \begin{equation}
    \label{eq:cross_attn_weights}
    \mA_{\text{cross}} = \begin{bmatrix}
    0.302 & 0.282 & 0.416 \\
    0.335 & 0.205 & 0.460
    \end{bmatrix} \in \R^{2 \times 3}.
    \end{equation}
    \item \textbf{Cross-attention context retrieval}:
    Multiplying $\mA_{\text{cross}} \in \R^{2 \times 3}$ by encoder values $\mV \in \R^{3 \times 2}$ yields the grounded representations $\mO_{\text{cross}} \in \R^{2 \times 2}$:
    \begin{equation}
    \label{eq:cross_output_matrix}
    \mO_{\text{cross}} = \mA_{\text{cross}} \mV = \begin{bmatrix}
    0.302 \vv_1 + 0.282 \vv_2 + 0.416 \vv_3 \\
    0.335 \vv_1 + 0.205 \vv_2 + 0.460 \vv_3
    \end{bmatrix}
    \approx \begin{bmatrix}
    0.648 & 1.117 \\
    0.706 & 1.151
    \end{bmatrix}.
    \end{equation}
\end{enumerate}

Query $\vq_{\text{``riva''}}$ strongly aligns with source key $\vk_{\text{``river''}}$ ($\alpha_{23} = 0.460$) and source key $\vk_{\text{``bank''}}$ ($\alpha_{21} = 0.335$). The decoder directly pulls the exact semantic coordinates of the geological shore from the source representation without compressing the source into a single bottleneck vector.

\clearpage
\section{Multi-Head Attention}
\label{sec:multi_head_attention}


In single-head attention, the model computes a single set of attention weights $\alpha_{i, j}$, averaging all values into a single weighted sum. However, a single token often participates in multiple simultaneous linguistic relationships:
\begin{itemize}
    \item A verb must track both its grammatical subject (often located earlier in the sentence) and its direct object (often located downstream).
    \item In programming languages, an identifier must link to its type declaration, its scope boundary, and its mutating assignments.
\end{itemize}
A single attention head cannot focus simultaneously on disparate positions with different semantic intents: the softmax distribution is forced to average across them, diluting specific relational signals.

\begin{figure}[htbp]
\centering
\resizebox{0.96\textwidth}{!}{%
\begin{tikzpicture}[
    >=Stealth,
    font=\small\sffamily,
    tokenbox/.style={rectangle, rounded corners=4pt, draw=politoNavy, fill=white, thick, minimum width=1.6cm, minimum height=0.7cm, align=center, font=\small\ttfamily},
    tokenbox_focus/.style={rectangle, rounded corners=4pt, draw=crimsonRed, fill=crimsonRed!15, very thick, minimum width=1.6cm, minimum height=0.7cm, align=center, font=\small\ttfamily\bfseries},
    panel/.style={rectangle, rounded corners=6pt, draw=lightGrayLine, fill=white, inner sep=10pt}
]

% Panel A: Single-Head Attention
\node[panel, minimum width=6.8cm, minimum height=5.2cm] (panel_a) at (0, 0) {};
\node[font=\bfseries\normalsize, text=politoNavy] at (0, 2.2) {(a) Single-Head Attention (Averaged / Diluted)};

\node[tokenbox] (a1) at (-2.4, 0.4) {The};
\node[tokenbox] (a2) at (-0.8, 0.4) {patient};
\node[tokenbox_focus] (a3) at (0.8, 0.4) {received};
\node[tokenbox] (a4) at (2.4, 0.4) {treatment};

% Blurry/averaged arrows from 'received'
\draw[->, line width=1.5pt, draw=bodyGray!80] (a3.north) to[bend right=35] node[above, font=\scriptsize, text=bodyDark] {$\alpha \approx 0.35$} (a2.north);
\draw[->, line width=1.5pt, draw=bodyGray!80] (a3.north) to[bend left=35] node[above, font=\scriptsize, text=bodyDark] {$\alpha \approx 0.45$} (a4.north);
\draw[->, line width=0.8pt, draw=bodyGray!50, dashed] (a3.north) to[bend right=45] node[above, font=\scriptsize, text=bodyGray] {$\alpha \approx 0.15$} (a1.north);

\node[align=center, font=\scriptsize, text=crimsonRed, fill=crimsonRed!5, rounded corners=3pt, draw=crimsonRed!30, inner sep=4pt] at (0, -1.5) {
    \textbf{Superposition Collapse}: Softmax mixes subject (\emph{patient})\\and direct object (\emph{treatment}) into a single averaged vector.\\Specific syntactic identities are mutually blurred.
};

% Panel B: Multi-Head Attention
\node[panel, minimum width=7.4cm, minimum height=5.2cm] (panel_b) at (8.0, 0) {};
\node[font=\bfseries\normalsize, text=politoNavy] at (8.0, 2.2) {(b) Multi-Head Attention (Specialized Subspaces)};

\node[tokenbox] (b1) at (5.6, 0.4) {The};
\node[tokenbox] (b2) at (7.2, 0.4) {patient};
\node[tokenbox_focus] (b3) at (8.8, 0.4) {received};
\node[tokenbox] (b4) at (10.4, 0.4) {treatment};

% Head 1: Subject relationship (Sky Blue)
\draw[->, very thick, draw=steelBlue] (b3.north) to[bend right=38] node[above=2pt, font=\scriptsize\bfseries, text=steelBlue] {Head 1: Subject ($\alpha = 0.88$)} (b2.north);

% Head 2: Direct Object relationship (Deepdive Teal)
\draw[->, very thick, draw=deepdiveTeal] (b3.south) to[bend right=38] node[below=2pt, font=\scriptsize\bfseries, text=deepdiveTeal] {Head 2: Direct Object ($\alpha = 0.92$)} (b4.south);

% Head 3: Determiner / Context (Exam Amber)
\draw[->, thick, draw=examAmber, dashed] (b2.south) to[bend left=30] node[below=2pt, font=\scriptsize\bfseries, text=examAmber] {Head 3: Specifier} (b1.south);

\node[align=center, font=\scriptsize, text=deepdiveTeal, fill=deepdiveTealBg, rounded corners=3pt, draw=deepdiveTealLine!50, inner sep=4pt] at (8.0, -1.5) {
    \textbf{Disentangled Linguistic Subspaces}:\\Each head linearly projects to independent $\R^{d_k}$ subspace.\\Head 1 isolates subject, Head 2 isolates direct object,\\preserving distinct syntactic roles simultaneously.
};

\end{tikzpicture}%
}
\caption{Linguistic relationship modeling: Single-Head vs. Multi-Head Attention. In Single-Head Attention (a), a token involved in multiple syntactic dependencies must compromise its single softmax distribution, averaging disparate features and diluting specific relational roles. In Multi-Head Attention (b), projection into $h$ distinct representation subspaces enables separate heads to concurrently isolate orthogonal linguistic relationships (such as grammatical subject tracking and direct object binding) without mutual interference.}
\label{fig:single_vs_multi_head_linguistic}
\end{figure}


Multi-Head Attention (MHA) resolves this by linearly projecting queries, keys, and values $h$ times with different, learned parameter matrices:
\begin{definition}{Multi-head attention}{mha_defn}
Let $h$ denote the number of attention heads. Multi-Head Attention projects $\mQ, \mK, \mV$ into $h$ distinct representation subspaces of dimension $d_k = \dmodel / h$ and $d_v = \dmodel / h$:
\begin{equation}
\label{eq:mha_formula}
\text{\texttt{MultiHead}}(\mQ, \mK, \mV) = \left[ \text{head}_1 \parallel \text{head}_2 \parallel \dots \parallel \text{head}_h \right] \mW^O,
\end{equation}
where each individual head is evaluated via:
\begin{equation}
\text{head}_i = \text{\texttt{Attention}}(\mQ \mW_i^Q, \mK \mW_i^K, \mV \mW_i^V), \quad i \in \{1, 2, \dots, h\},
\end{equation}
with parameter projection matrices $\mW_i^Q \in \R^{\dmodel \times d_k}$, $\mW_i^K \in \R^{\dmodel \times d_k}$, $\mW_i^V \in \R^{\dmodel \times d_v}$, and output projection matrix $\mW^O \in \R^{h d_v \times \dmodel}$.
\end{definition}

\begin{figure}[htbp]
\centering
\begin{tikzpicture}[
    scale=0.88, every node/.style={transform shape},
    box/.style={draw=politoNavy, fill=skyBlue!20, thick, rounded corners=3pt, minimum width=1.6cm, minimum height=0.65cm, align=center, font=\footnotesize\sffamily},
    attnhead/.style={draw=examAmberLine, fill=examAmberBg, thick, rounded corners=3pt, minimum width=1.8cm, minimum height=0.7cm, align=center, font=\footnotesize\sffamily\bfseries},
    arrow/.style={-{Stealth[length=5pt]}, thick, draw=politoNavy}
]
    \node[font=\sffamily\bfseries] (q_in) at (-3.2, 0) {$\mQ \in \R^{L \times \dmodel}$};
    \node[font=\sffamily\bfseries] (k_in) at (0, 0) {$\mK \in \R^{L \times \dmodel}$};
    \node[font=\sffamily\bfseries] (v_in) at (3.2, 0) {$\mV \in \R^{L \times \dmodel}$};

    % LINEAR PROJECTIONS
    \node[box] (l_q1) at (-4.2, 1.3) {Linear 1};
    \node[box] (l_q2) at (-2.2, 1.3) {Linear $h$};
    \node at (-3.2, 1.3) {$\dots$};

    \node[box] (l_k1) at (-1.0, 1.3) {Linear 1};
    \node[box] (l_k2) at (1.0, 1.3) {Linear $h$};
    \node at (0, 1.3) {$\dots$};

    \node[box] (l_v1) at (2.2, 1.3) {Linear 1};
    \node[box] (l_v2) at (4.2, 1.3) {Linear $h$};
    \node at (3.2, 1.3) {$\dots$};

    \draw[arrow] (q_in) -- (l_q1); \draw[arrow] (q_in) -- (l_q2);
    \draw[arrow] (k_in) -- (l_k1); \draw[arrow] (k_in) -- (l_k2);
    \draw[arrow] (v_in) -- (l_v1); \draw[arrow] (v_in) -- (l_v2);

    % ATTENTION HEADS
    \node[attnhead] (head1) at (-2.6, 2.7) {Scaled Dot-Prod \\ Head 1};
    \node[attnhead] (headh) at (2.6, 2.7) {Scaled Dot-Prod \\ Head $h$};
    \node at (0, 2.7) {$\dots$};

    \draw[arrow] (l_q1) -- (head1); \draw[arrow] (l_k1) -- (head1); \draw[arrow] (l_v1) -- (head1);
    \draw[arrow] (l_q2) -- (headh); \draw[arrow] (l_k2) -- (headh); \draw[arrow] (l_v2) -- (headh);

    % CONCATENATION & PROJECTION
    \node[box, fill=politoNavy!15, minimum width=7.5cm, minimum height=0.7cm] (concat) at (0, 4.0) {Concat $[\text{head}_1 \parallel \dots \parallel \text{head}_h] \in \R^{L \times h d_v}$};
    \draw[arrow] (head1) -- (concat.south -| head1);
    \draw[arrow] (headh) -- (concat.south -| headh);

    \node[box, fill=thmGreenBg, draw=thmGreenLine, minimum width=5.0cm] (proj) at (0, 5.2) {Linear Projection $\mW^O \in \R^{h d_v \times \dmodel}$};
    \draw[arrow] (concat) -- (proj);

    \node[above=0.6cm of proj, font=\sffamily\bfseries\color{politoNavy}] (mha_out) {Output $\in \R^{L \times \dmodel}$};
    \draw[arrow] (proj) -- (mha_out);
\end{tikzpicture}
\caption{Architecture of Multi-Head Attention (MHA) \cite{vaswani2017}. Projections map inputs into $h$ parallel representation subspaces of reduced dimension $d_k = d_v = \dmodel / h$, whose outputs are concatenated and linearly projected back to $\dmodel$.}
\label{fig:multi_head_attention_tikz}
\end{figure}

\subsubsection{Parameter equivalence and computational efficiency}
Due to the reduced dimension of each head ($d_k = d_v = \dmodel / h$), the total computational cost of multi-head attention is identical to that of single-head attention with full dimensionality $\dmodel$:
\begin{equation}
h \times \left( \dmodel \times \frac{\dmodel}{h} \right) = \dmodel^2 \text{ operations for projections}.
\end{equation}
In standard configurations (e.g., $\dmodel = 512, h = 8 \implies d_k = d_v = 64$, or $\dmodel = 768, h = 12 \implies d_k = d_v = 64$), parallel heads learn specialized sub-behaviors: some heads attend to local syntactic structure, others to long-range dependencies, and others to delimiter tokens.

\clearpage
\section{Residual Connections and Normalization}
\label{sec:sublayer_mechanics}

In deep Transformer architectures comprising dozens of layers, stacking attention blocks and non-linear feed-forward networks directly would lead to optimization failure due to exploding or vanishing gradients. To ensure stable gradient propagation, each Transformer sublayer is wrapped in a residual connection followed by Layer Normalization:
\begin{equation}
\label{eq:add_and_norm_post}
\vy = \text{\texttt{LayerNorm}}(\vx + \text{\texttt{Sublayer}}(\vx)).
\end{equation}

\subsubsection{Residual skip connections}
First introduced in \cite{he2016deep} for deep residual networks (ResNets), identity shortcut connections pass the unmodified input $\vx$ around the sublayer (Figure~\ref{fig:residual_connections}):
\begin{equation}
\label{eq:residual_connection}
\vy = \vx + \mathcal{F}(\vx),
\end{equation}
where $\mathcal{F}(\vx)$ represents the multi-head attention or feed-forward transformation.

\begin{figure}[htbp]
\centering
\resizebox{0.75\textwidth}{!}{%
\begin{tikzpicture}[
    >=Stealth,
    font=\small\sffamily,
    box/.style={rectangle, rounded corners=5pt, draw=politoNavy, fill=skyBlue!25, thick, minimum width=4.5cm, minimum height=1.4cm, align=center},
    addnode/.style={circle, draw=politoNavy, fill=white, very thick, minimum size=8mm, font=\large\bfseries, inner sep=0pt},
    arr/.style={-{Stealth[scale=1.1]}, thick, draw=politoNavy},
    skiparr/.style={-{Stealth[scale=1.1]}, very thick, draw=crimsonRed},
    gradarr/.style={-{Stealth[scale=1.1]}, thick, dashed, draw=crimsonRed}
]

% Forward diagram
\node (x) at (0, 0) {$\vx \in \R^{\dmodel}$};
\coordinate (split) at (0, 1.0);
\node[box] (sublayer) at (0, 2.5) {\textbf{Transformer Sublayer}\\$\mathcal{F}(\vx)$ (MHA or FFN)};
\node[addnode] (add) at (0, 4.2) {+};
\node (y) at (0, 5.4) {$\vy = \vx + \mathcal{F}(\vx)$};

% Forward connections
\draw[arr] (x) -- (split) -- (sublayer);
\draw[arr] (sublayer) -- (add);
\draw[arr] (add) -- (y);
\draw[skiparr] (split) -- ++(-2.4, 0) -- node[left, font=\footnotesize\bfseries, text=crimsonRed] {Identity skip connection ($\vx$)} ++(0, 3.2) -- (add);

% Annotations & gradient equations
\node[draw=lightGrayLine, rounded corners=4pt, fill=examAmberBg, inner sep=8pt, align=left, font=\footnotesize] at (4.6, 2.5) {
    \textbf{Gradient superhighway}:\\[3pt]
    $\displaystyle \frac{\partial \mathcal{L}}{\partial \vx} = \frac{\partial \mathcal{L}}{\partial \vy} \left( \mI + \frac{\partial \mathcal{F}(\vx)}{\partial \vx} \right)$\\[6pt]
    Even if $\displaystyle \frac{\partial \mathcal{F}}{\partial \vx} \to \mat{0}$, the constant identity\\[2pt]
    term $\mI$ propagates gradients directly to\\
    earlier layers without attenuation.
};

\end{tikzpicture}%
}
\caption{Residual skip connection and gradient propagation dynamics. The input $\mathbf{x}$ bypasses the non-linear sublayer $\mathcal{F}(\mathbf{x})$ via an identity shortcut. During backpropagation, the derivative $\frac{\partial \mathbf{y}}{\partial \mathbf{x}} = \mathbf{I} + \frac{\partial \mathcal{F}}{\partial \mathbf{x}}$ ensures that the gradient $\frac{\partial \mathcal{L}}{\partial \mathbf{y}}$ flows unimpeded across arbitrary layer depths via the identity term $\mathbf{I}$.}
\label{fig:residual_connections}
\end{figure}

During backpropagation, the gradient of the loss $\mathcal{L}$ with respect to input $\vx$ expands via the chain rule to:
\begin{equation}
\label{eq:residual_gradient_flow}
\frac{\partial \mathcal{L}}{\partial \vx} = \frac{\partial \mathcal{L}}{\partial \vy} \frac{\partial \vy}{\partial \vx} = \frac{\partial \mathcal{L}}{\partial \vy} \left( \mI + \frac{\partial \mathcal{F}(\vx)}{\partial \vx} \right) = \frac{\partial \mathcal{L}}{\partial \vy} + \frac{\partial \mathcal{L}}{\partial \vy} \frac{\partial \mathcal{F}(\vx)}{\partial \vx}.
\end{equation}
The presence of the identity matrix $\mI$ guarantees an unobstructed \textbf{gradient superhighway}: even if the Jacobian $\frac{\partial \mathcal{F}}{\partial \vx}$ vanishes (for instance, when weights are initialized small or activations saturate), term $\frac{\partial \mathcal{L}}{\partial \vy}$ flows directly back to earlier layers in a single step, completely preventing gradient extinction.

\subsubsection{Layer Normalization}
Unlike Batch Normalization (which computes statistics across the mini-batch dimension, making it sensitive to batch size and variable sequence lengths), \textbf{Layer Normalization} \cite{ba2016layer} computes normalization statistics across the feature dimensions for each individual token independently (Figure~\ref{fig:layer_normalization}):

\begin{figure}[htbp]
\centering
\resizebox{0.92\textwidth}{!}{%
\begin{tikzpicture}[
    >=Stealth,
    font=\small\sffamily,
    panel/.style={draw=lightGrayLine, rounded corners=6pt, fill=white, inner sep=10pt},
    header/.style={font=\bfseries\small, align=center},
    lbl/.style={font=\scriptsize, align=center, text=bodyDark}
]

% Left: Batch Normalization (BN)
\node[panel, minimum width=6.2cm, minimum height=6.2cm] (p1) at (0,0) {};
\node[header, text=crimsonRed] at (0, 2.5) {Batch Normalization (BN)\\\textmd{\footnotesize Across batch samples for each feature}};

% 3D cube for BN
\draw[thick, draw=bodyGray] (-2, -1.2) rectangle (1.2, 1.2);
\draw[thick, draw=bodyGray] (1.2, -1.2) -- ++(0.9, 0.7) -- ++(0, 2.4) -- (1.2, 1.2);
\draw[thick, draw=bodyGray] (-2, 1.2) -- ++(0.9, 0.7) -- ++(3.2, 0) -- (1.2, 1.2);

% Slice for BN: vertical column along Batch dimension
\fill[crimsonRed!40, draw=crimsonRed, thick, opacity=0.85] (-0.6, -1.2) rectangle (-0.1, 1.2);
\fill[crimsonRed!30, draw=crimsonRed, thick, opacity=0.85] (-0.1, -1.2) -- ++(0.9, 0.7) -- ++(0, 2.4) -- (-0.1, 1.2) -- cycle;
\fill[crimsonRed!20, draw=crimsonRed, thick, opacity=0.85] (-0.6, 1.2) -- ++(0.9, 0.7) -- ++(0.5, 0) -- (-0.1, 1.2) -- cycle;

\node[font=\tiny\bfseries, text=crimsonRed] at (-0.35, 0) {Normalize};
\node[lbl] at (0, -1.8) {$\displaystyle \mu_d = \frac{1}{B} \sum_{b=1}^B x_{b,t,d}$};
\node[lbl, text=crimsonRed] at (0, -2.4) {\textbf{NLP Failure Mode}: Dependent on batch size\\and disrupted by variable-length padding};

% Right: Layer Normalization (LN)
\node[panel, minimum width=6.2cm, minimum height=6.2cm] (p2) at (7.2,0) {};
\node[header, text=steelBlue] at (7.2, 2.5) {Layer Normalization (LN)\\\textmd{\footnotesize Across features for each individual token}};

% 3D cube for LN
\draw[thick, draw=bodyGray] (5.2, -1.2) rectangle (8.4, 1.2);
\draw[thick, draw=bodyGray] (8.4, -1.2) -- ++(0.9, 0.7) -- ++(0, 2.4) -- (8.4, 1.2);
\draw[thick, draw=bodyGray] (5.2, 1.2) -- ++(0.9, 0.7) -- ++(3.2, 0) -- (8.4, 1.2);

% Slice for LN: horizontal slice across all features
\fill[steelBlue!40, draw=steelBlue, thick, opacity=0.85] (5.2, 0.2) rectangle (8.4, 0.6);
\fill[steelBlue!30, draw=steelBlue, thick, opacity=0.85] (8.4, 0.2) -- ++(0.9, 0.7) -- ++(0, 0.4) -- (8.4, 0.6) -- cycle;
\fill[steelBlue!20, draw=steelBlue, thick, opacity=0.85] (5.2, 0.6) -- ++(0.9, 0.7) -- ++(3.2, 0) -- (8.4, 0.6) -- cycle;

\node[font=\tiny\bfseries, text=politoNavy] at (6.8, 0.4) {Normalize};
\node[lbl] at (7.2, -1.8) {$\displaystyle \mu_{b,t} = \frac{1}{D} \sum_{d=1}^D x_{b,t,d}$};
\node[lbl, text=steelBlue] at (7.2, -2.4) {\textbf{NLP Standard}: Invariant to batch size,\\completely independent across time steps};

\end{tikzpicture}%
}
\caption{Comparison of normalization geometries in 3D activation tensor space $\mathcal{X} \in \mathbb{R}^{B \times T \times D}$. (Left) Batch Normalization computes statistics across all batch instances $B$ for each feature channel $d$, creating inter-sample dependencies and breaking down under variable sequence lengths and sequence padding. (Right) Layer Normalization computes statistics across all feature channels $D$ for each token vector $\mathbf{x}_{b,t}$ independently, ensuring strict batch-size invariance and identical execution during training and autoregressive inference.}
\label{fig:layer_normalization}
\end{figure}

\begin{definition}{Layer Normalization}{layer_norm_defn}
For a token representation $\vx = [x_1, x_2, \dots, x_d]^\top \in \R^d$, Layer Normalization evaluates:
\begin{align}
\mu &= \frac{1}{d} \sum_{j=1}^d x_j, \quad \sigma^2 = \frac{1}{d} \sum_{j=1}^d (x_j - \mu)^2, \\
\label{eq:layer_norm_formula}
\text{\texttt{LayerNorm}}(\vx) &= \frac{\vx - \mu}{\sqrt{\sigma^2 + \epsilon}} \odot \boldsymbol{\gamma} + \boldsymbol{\beta},
\end{align}
where $\epsilon > 0$ is a small numerical stabilization constant, and $\boldsymbol{\gamma}, \boldsymbol{\beta} \in \R^d$ are learnable affine gain and bias parameters initialized to $\boldsymbol{\gamma} = \mathbf{1}$ and $\boldsymbol{\beta} = \mathbf{0}$.
\end{definition}
Because Layer Normalization operates on a per-token basis without batch-level pooling, its behavior is identical during training and evaluation, providing exceptional stability across variable-length sequences.

\subsubsection{Position-wise Feed-Forward Networks (FFN)}
In addition to attention sublayers, each layer in the encoder and decoder contains a fully connected \textbf{Feed-Forward Network} applied to each position separately and identically:
\begin{equation}
\label{eq:position_wise_ffn}
\text{\texttt{FFN}}(\vx) = \max\left(0, \vx \mW_1 + \vb_1\right) \mW_2 + \vb_2,
\end{equation}
where $\mW_1 \in \R^{\dmodel \times \dff}$, $\vb_1 \in \R^{\dff}$, $\mW_2 \in \R^{\dff \times \dmodel}$, and $\vb_2 \in \R^{\dmodel}$. Typically, the inner hidden dimension is expanded by a factor of 4 ($\dff = 4\dmodel$, e.g., $\dff = 2048$ for $\dmodel = 512$). While attention routes information across different temporal positions, the FFN provides non-linear feature transformation within each individual position.

\subsection{End-to-end sublayer integration: the complete mathematical pipeline}
\label{subsec:end_to_end_sublayer_pipeline}

Having analyzed Multi-Head Attention, residual skip connections, Layer Normalization, and Feed-Forward Networks individually, we now formalize how these mathematical operators chain together within complete Transformer layers, tracing their exact locations in the canonical architecture (Figure~\ref{fig:vaswani_highlight_sublayers}).

\begin{figure}[htbp]
\centering
\begin{subfigure}[b]{0.47\textwidth}
    \centering
    \resizebox{\linewidth}{!}{%
    \begin{tikzpicture}[
        >=Stealth,
        font=\small\sffamily,
        block/.style={rectangle, draw=tfOrangeBorder, fill=tfOrange, thick, rounded corners=4pt, minimum width=3.4cm, minimum height=0.7cm, align=center, font=\footnotesize\bfseries},
        normblock/.style={rectangle, draw=tfYellowBorder, fill=tfYellow, thick, rounded corners=4pt, minimum width=3.4cm, minimum height=0.6cm, align=center, font=\footnotesize\bfseries},
        ffnblock/.style={rectangle, draw=tfBlueBorder, fill=tfBlue, thick, rounded corners=4pt, minimum width=3.4cm, minimum height=0.7cm, align=center, font=\footnotesize\bfseries},
        sum/.style={circle, draw=politoNavy, fill=white, thick, minimum size=0.5cm, font=\bfseries\small, inner sep=0pt},
        box/.style={draw=politoNavy, dashed, fill=tfGrayBg, rounded corners=6pt, inner sep=8pt}
    ]
        % Outer layer container (N x)
        \node[box, minimum width=4.8cm, minimum height=6.3cm] (encstack) at (0, 3.4) {};
        \node[font=\bfseries\footnotesize, text=politoNavy, anchor=north east] at (encstack.north east) [xshift=-4pt, yshift=-4pt] {$N\times$};
        \node[font=\footnotesize\bfseries\scshape, text=politoNavy, anchor=north west] at (encstack.north west) [xshift=6pt, yshift=-4pt] {Encoder Layer $l$};

        % Input
        \node (in) at (0, 0) {$\mH_{\text{enc}}^{(l-1)} \in \R^{L \times \dmodel}$};
        \coordinate (split1) at (0, 0.7);
        \draw[thick] (in) -- (split1);

        % Multi-Head Attention
        \node[block] (mha) at (0, 1.4) {Multi-Head Attention};
        \draw[->, thick] (split1) -- (mha.south);
        \draw[->, thick] (split1) -| ++(-0.9, 0.3) -- ++(0, 0.4);
        \draw[->, thick] (split1) -| ++(0.9, 0.3) -- ++(0, 0.4);

        % Add 1
        \node[sum] (add1) at (0, 2.3) {$\oplus$};
        \draw[->, thick] (mha) -- (add1);

        % Residual 1
        \draw[->, thick, rounded corners=4pt, draw=crimsonRed] (split1) -- ++(-2.0, 0) -- ++(0, 1.6) -- (add1.west);

        % LayerNorm 1
        \node[normblock] (ln1) at (0, 3.1) {Add \& Norm (LN)};
        \draw[->, thick] (add1) -- (ln1);

        % Split 2
        \coordinate (split2) at (0, 3.7);
        \draw[thick] (ln1) -- (split2);

        % Feed Forward
        \node[ffnblock] (ffn) at (0, 4.4) {Feed-Forward (FFN)};
        \draw[->, thick] (split2) -- (ffn);

        % Add 2
        \node[sum] (add2) at (0, 5.3) {$\oplus$};
        \draw[->, thick] (ffn) -- (add2);

        % Residual 2
        \draw[->, thick, rounded corners=4pt, draw=crimsonRed] (split2) -- ++(-2.0, 0) -- ++(0, 1.6) -- (add2.west);

        % LayerNorm 2
        \node[normblock] (ln2) at (0, 6.0) {Add \& Norm (LN)};
        \draw[->, thick] (add2) -- (ln2);

        % Output
        \node (out) at (0, 7.0) {$\mH_{\text{enc}}^{(l)} \in \R^{L \times \dmodel}$};
        \draw[->, very thick] (ln2) -- (out);

    \end{tikzpicture}%
    }
    \caption{Encoder layer stack ($N\times$)}
    \label{fig:vaswani_enc_sublayers}
\end{subfigure}\hfill
\begin{subfigure}[b]{0.51\textwidth}
    \centering
    \resizebox{\linewidth}{!}{%
    \begin{tikzpicture}[
        >=Stealth,
        font=\small\sffamily,
        block/.style={rectangle, draw=tfOrangeBorder, fill=tfOrange, thick, rounded corners=4pt, minimum width=3.4cm, minimum height=0.65cm, align=center, font=\footnotesize\bfseries},
        normblock/.style={rectangle, draw=tfYellowBorder, fill=tfYellow, thick, rounded corners=4pt, minimum width=3.4cm, minimum height=0.55cm, align=center, font=\footnotesize\bfseries},
        ffnblock/.style={rectangle, draw=tfBlueBorder, fill=tfBlue, thick, rounded corners=4pt, minimum width=3.4cm, minimum height=0.65cm, align=center, font=\footnotesize\bfseries},
        sum/.style={circle, draw=politoNavy, fill=white, thick, minimum size=0.45cm, font=\bfseries\small, inner sep=0pt},
        box/.style={draw=forestGreen, dashed, fill=thmGreenBg!30, rounded corners=6pt, inner sep=8pt}
    ]
        % Outer layer container (N x)
        \node[box, minimum width=5.2cm, minimum height=8.6cm] (decstack) at (0, 4.5) {};
        \node[font=\bfseries\footnotesize, text=forestGreen, anchor=north east] at (decstack.north east) [xshift=-4pt, yshift=-4pt] {$N\times$};
        \node[font=\footnotesize\bfseries\scshape, text=forestGreen, anchor=north west] at (decstack.north west) [xshift=6pt, yshift=-4pt] {Decoder Layer $l$};

        % Input
        \node (in) at (0, 0) {$\mH_{\text{dec}}^{(l-1)} \in \R^{L_{\text{dec}} \times \dmodel}$};
        \coordinate (split1) at (0, 0.6);
        \draw[thick] (in) -- (split1);

        % Sublayer 1: Masked MHA
        \node[block] (mmha) at (0, 1.25) {Masked Multi-Head Attn};
        \draw[->, thick] (split1) -- (mmha.south);
        \draw[->, thick] (split1) -| ++(-0.8, 0.25) -- ++(0, 0.4);
        \draw[->, thick] (split1) -| ++(0.8, 0.25) -- ++(0, 0.4);

        \node[sum] (add1) at (0, 2.05) {$\oplus$};
        \draw[->, thick] (mmha) -- (add1);
        \draw[->, thick, rounded corners=4pt, draw=crimsonRed] (split1) -- ++(-2.0, 0) -- ++(0, 1.45) -- (add1.west);

        \node[normblock] (ln1) at (0, 2.75) {Add \& Norm (LN)};
        \draw[->, thick] (add1) -- (ln1);

        % Sublayer 2: Cross MHA
        \coordinate (split2) at (0, 3.4);
        \draw[thick] (ln1) -- (split2);

        \node[block] (cmha) at (0, 4.15) {Cross Multi-Head Attn};
        \draw[->, thick] (split2) -- node[right, font=\tiny]{$\mQ$} (cmha.south);

        % External Encoder Input
        \node[anchor=east, font=\scriptsize\bfseries, text=deepdiveTeal] (encin) at (-1.7, 4.15) {$\mH_{\text{enc}}^{(N)}$};
        \draw[->, thick, draw=deepdiveTeal] (encin.east) -- node[above, font=\tiny]{$\mK,\mV$} (cmha.west);

        \node[sum] (add2) at (0, 5.0) {$\oplus$};
        \draw[->, thick] (cmha) -- (add2);
        \draw[->, thick, rounded corners=4pt, draw=crimsonRed] (split2) -- ++(2.0, 0) -- ++(0, 1.6) -- (add2.east);

        \node[normblock] (ln2) at (0, 5.75) {Add \& Norm (LN)};
        \draw[->, thick] (add2) -- (ln2);

        % Sublayer 3: FFN
        \coordinate (split3) at (0, 6.4);
        \draw[thick] (ln2) -- (split3);

        \node[ffnblock] (ffn) at (0, 7.15) {Feed-Forward (FFN)};
        \draw[->, thick] (split3) -- (ffn);

        \node[sum] (add3) at (0, 8.0) {$\oplus$};
        \draw[->, thick] (ffn) -- (add3);
        \draw[->, thick, rounded corners=4pt, draw=crimsonRed] (split3) -- ++(-2.0, 0) -- ++(0, 1.6) -- (add3.west);

        \node[normblock] (ln3) at (0, 8.75) {Add \& Norm (LN)};
        \draw[->, thick] (add3) -- (ln3);

        % Output
        \node (out) at (0, 9.6) {$\mH_{\text{dec}}^{(l)} \in \R^{L_{\text{dec}} \times \dmodel}$};
        \draw[->, very thick] (ln3) -- (out);

    \end{tikzpicture}%
    }
    \caption{Decoder layer stack ($N\times$)}
    \label{fig:vaswani_dec_sublayers}
\end{subfigure}
\caption{End-to-end sublayer connectivity within the canonical Transformer architecture \cite{vaswani2017}. Every attention and feed-forward block is coupled with an identity residual shortcut and an ``Add \& Norm'' operator.}
\label{fig:vaswani_highlight_sublayers}
\end{figure}

\subsubsection{Complete Encoder layer pipeline}
An encoder layer $l \in \{1, \dots, N\}$ receives tensor $\mH_{\text{enc}}^{(l-1)} \in \R^{L_{\text{enc}} \times \dmodel}$ (with $\mH_{\text{enc}}^{(0)} = \mX_{\text{enc}}$) and applies two composite sublayers:
\begin{enumerate}
    \item \textbf{Sublayer 1 (Bidirectional Multi-Head Self-Attention)}:
    \begin{align}
    \label{eq:enc_sublayer1_attn}
    \mZ^{(l)} &= \text{\texttt{MultiHead}}\left(\mH_{\text{enc}}^{(l-1)}, \mH_{\text{enc}}^{(l-1)}, \mH_{\text{enc}}^{(l-1)}\right), \\
    \label{eq:enc_sublayer1_norm}
    \tilde{\mH}_{\text{enc}}^{(l)} &= \text{\texttt{LayerNorm}}\left(\mH_{\text{enc}}^{(l-1)} + \mZ^{(l)}\right).
    \end{align}
    \item \textbf{Sublayer 2 (Position-wise Feed-Forward Network)}:
    \begin{align}
    \label{eq:enc_sublayer2_ffn}
    \mathbf{F}^{(l)} &= \text{\texttt{FFN}}\left(\tilde{\mH}_{\text{enc}}^{(l)}\right), \\
    \label{eq:enc_sublayer2_norm}
    \mH_{\text{enc}}^{(l)} &= \text{\texttt{LayerNorm}}\left(\tilde{\mH}_{\text{enc}}^{(l)} + \mathbf{F}^{(l)}\right).
    \end{align}
\end{enumerate}

\subsubsection{Complete Decoder layer pipeline}
A decoder layer $l \in \{1, \dots, N\}$ receives target representation $\mH_{\text{dec}}^{(l-1)} \in \R^{L_{\text{dec}} \times \dmodel}$ along with the terminal encoder output $\mH_{\text{enc}}^{(N)} \in \R^{L_{\text{enc}} \times \dmodel}$, executing three consecutive sublayers:
\begin{enumerate}
    \item \textbf{Sublayer 1 (Causal Masked Multi-Head Self-Attention)}:
    \begin{align}
    \label{eq:dec_sublayer1_attn}
    \mZ_{\text{self}}^{(l)} &= \text{\texttt{MaskedMultiHead}}\left(\mH_{\text{dec}}^{(l-1)}, \mH_{\text{dec}}^{(l-1)}, \mH_{\text{dec}}^{(l-1)}\right), \\
    \label{eq:dec_sublayer1_norm}
    \mH_{\text{dec}, 1}^{(l)} &= \text{\texttt{LayerNorm}}\left(\mH_{\text{dec}}^{(l-1)} + \mZ_{\text{self}}^{(l)}\right).
    \end{align}
    \item \textbf{Sublayer 2 (Encoder-Decoder Multi-Head Cross-Attention)}:
    \begin{align}
    \label{eq:dec_sublayer2_attn}
    \mZ_{\text{cross}}^{(l)} &= \text{\texttt{MultiHead}}\left(\mH_{\text{dec}, 1}^{(l)}, \mH_{\text{enc}}^{(N)}, \mH_{\text{enc}}^{(N)}\right), \\
    \label{eq:dec_sublayer2_norm}
    \mH_{\text{dec}, 2}^{(l)} &= \text{\texttt{LayerNorm}}\left(\mH_{\text{dec}, 1}^{(l)} + \mZ_{\text{cross}}^{(l)}\right).
    \end{align}
    Here, queries are projected from $\mH_{\text{dec}, 1}^{(l)}$, while keys and values are projected from $\mH_{\text{enc}}^{(N)}$.
    \item \textbf{Sublayer 3 (Position-wise Feed-Forward Network)}:
    \begin{align}
    \label{eq:dec_sublayer3_ffn}
    \mathbf{F}_{\text{dec}}^{(l)} &= \text{\texttt{FFN}}\left(\mH_{\text{dec}, 2}^{(l)}\right), \\
    \label{eq:dec_sublayer3_norm}
    \mH_{\text{dec}}^{(l)} &= \text{\texttt{LayerNorm}}\left(\mH_{\text{dec}, 2}^{(l)} + \mathbf{F}_{\text{dec}}^{(l)}\right).
    \end{align}
\end{enumerate}

\subsubsection{Post-LN vs. Pre-LN parameter stabilization}
The formulation above represents the original \textbf{Post-LN} architecture of \cite{vaswani2017}, where normalization is placed on the residual sum: $\vy = \text{\texttt{LN}}(\vx + \mathcal{F}(\vx))$.
While mathematically elegant, deep Post-LN models require delicate learning rate warm-up schedules because gradients passing through the unnormalized residual additions can destabilize early training.
Modern implementations (such as GPT-2, LLaMA, and Mistral) predominantly adopt the \textbf{Pre-LN} formulation:
\begin{equation}
\label{eq:pre_ln_formula}
\vy = \vx + \mathcal{F}\left(\text{\texttt{LayerNorm}}(\vx)\right),
\end{equation}
which preserves an unbroken, completely unnormalized identity shortcut throughout the entire depth of the network, enabling stable optimization without learning rate warm-up.

\clearpage
\section{Positional Representations}
\label{sec:relative_positional_embeddings}

In Chapter~\ref{ch:rnn_sequence_modeling}, we introduced absolute positional encodings---both fixed sinusoidal functions \cite{vaswani2017} and learned absolute embedding vectors. Under absolute encoding, each index in the sequence ($pos = 0, 1, 2, \dots$) is assigned a dedicated vector added directly to the token embedding before entering the Transformer.


In natural language and software engineering source code, the precise absolute coordinate of a token is frequently less important than its \textbf{relative displacement} with respect to other tokens. For example, whether a verb appears at index 12 or index 84, its grammatical subject is typically found 2 or 3 tokens earlier.
To incorporate this inductive bias directly into the attention mechanism, Shaw et al. \cite{shaw2018} formulated \textbf{Relative Positional Representations}.


Instead of adding absolute vectors to the input, relative positional embeddings inject pairwise distance information directly into the scaled dot-product attention logits :
\begin{equation}
\label{eq:relative_position_attention}
\text{\texttt{Attention}}(\mQ, \mK, \mV) = \softmax\left( \frac{\mQ \mK^\top + \mS_{\text{rel}}}{\sqrt{d_k}} \right) \mV,
\end{equation}
where $\mS_{\text{rel}} \in \R^{L \times L}$ is a relative position bias matrix. For query position $i$ and key position $j$, the entry is defined by:
\begin{equation}
S_{\text{rel}}[i, j] = \vq_i^\top \mathbf{a}_{\text{clip}(i - j, k)}^K,
\end{equation}
where $\mathbf{a}^K$ is a learned embedding representing the relative offset $i - j$, and $\text{clip}(x, k) = \max(-k, \min(k, x))$ truncates relative distances beyond a maximum clipping horizon $k$.

\subsubsection{Asymmetry of relative positional embeddings}
Relative positional embeddings make sense exclusively for self-attention, and cannot be naturally defined for cross-attention:
\begin{itemize}
    \item \textbf{Self-attention}: Both queries and keys originate from the same sequence. The scalar distance $i - j$ represents a physical, well-defined offset within that sequence (e.g., token $j$ occurs 2 steps before token $i$).
    \item \textbf{Cross-attention}: Queries stem from the target sequence (e.g., generated English words in machine translation), while keys stem from the source sequence (e.g., input Italian words). Because source and target sentences possess different lengths, word orders, and grammatical alignments, there is no shared coordinate frame: asking whether an English target word is ``three positions ahead'' of an Italian source word is mathematically meaningless. Consequently, cross-attention cannot leverage relative positional biases.
\end{itemize}

\clearpage
\section{Computational Complexity}
\label{sec:computational_complexity}

To appreciate why the Transformer superseded Recurrent Neural Networks, we analyze three foundational computational metrics: computational complexity per layer, sequential operations, and maximum signal path length \cite{vaswani2017}, summarized in Table~\ref{tab:layer_complexity_comparison}.

\begin{table}[htbp]
\centering
\small
\caption{Computational complexity and operational parallelism across sequence modeling layers \cite{vaswani2017}.}
\label{tab:layer_complexity_comparison}
\begin{tabularx}{\textwidth}{lXXX}
\toprule
\textbf{Layer type} & \textbf{Complexity per layer} & \textbf{Sequential operations} & \textbf{Maximum path length} \\
\midrule
\textbf{Self-Attention} & $\mathcal{O}(L^2 \cdot d)$ & $\mathcal{O}(1)$ & $\mathcal{O}(1)$ \\
\textbf{Recurrent (RNN/LSTM)} & $\mathcal{O}(L \cdot d^2)$ & $\mathcal{O}(L)$ & $\mathcal{O}(L)$ \\
\textbf{Convolutional} & $\mathcal{O}(k \cdot L \cdot d^2)$ & $\mathcal{O}(1)$ & $\mathcal{O}(\log_k(L))$ \\
\bottomrule
\end{tabularx}
\end{table}

\subsubsection{Sequential operations: shattering the recurrence barrier}
In an RNN, computing the representation at step $L$ requires traversing an unbroken chain of $L$ sequential state updates ($\mathcal{O}(L)$ sequential operations).
In Self-Attention, computing all pairwise dot products $\mQ \mK^\top$ and value weighted sums involves dense matrix multiplications that execute in a single parallel step ($\mathcal{O}(1)$ sequential operations) across GPU tensor cores.

\subsubsection{Maximum path length: constant-time gradient propagation}
In recurrent networks, information and backpropagation gradients traveling between token $1$ and token $L$ must traverse $L$ intermediate state transitions. As proven by the Jacobian product $\prod_{k=t+1}^T \frac{\partial \vh_k}{\partial \vh_{k-1}}$ in Chapter~\ref{ch:rnn_sequence_modeling}, long path lengths cause exponential gradient decay.
In Self-Attention, the distance between any arbitrary pair of tokens $(i, j)$ is strictly $\mathcal{O}(1)$: the attention weight $\alpha_{i, j}$ directly links token $i$ to token $j$ in a single operation. This constant path length completely eliminates topological distance as an obstacle to learning long-range dependencies.

\subsubsection{Asymptotic trade-off: sequence length vs. latent dimension}
Self-attention incurs a quadratic time and memory complexity $\mathcal{O}(L^2 \cdot d)$ with respect to sequence length $L$.
When sequence length $L$ is smaller than embedding dimension $d$ (as was typical in early machine translation benchmarks where $L \approx 50$ and $d = 512$), $\mathcal{O}(L^2 \cdot d)$ is strictly faster than the recurrent cost $\mathcal{O}(L \cdot d^2)$. While long context windows ($L > 8192$) eventually encounter the $\mathcal{O}(L^2)$ memory wall, the massive parallel speedup during training outweighs this cost for standard operational contexts.

\clearpage

\clearpage
\thispagestyle{empty}
\null
\clearpage

\end{document}